% Alternate 1_2, derived directly from the explicitly supplied source pair.
% Build from this directory; all dependencies remain in this academic-year tree.
\makeatletter
\def\input@path{{../../../preamble/}}
\makeatother
\input{preamble_exams.tex}
\input{highlights/highlights.tex}
% \GTenHighlightsfalse % Uncomment to disable instructional highlights.
\input{highlights/specific_highlights.tex}
\CoordinatedTermTwoPalette
\tikzset{gten marker/.append style={inner xsep=.6mm,inner ysep=.3mm}}
\usepackage{varwidth}
\usepackage{needspace}
% Natural cell widths, bounded locally so long descriptions can wrap.
\newcommand{\CreditDescription}[1]{\begin{varwidth}[t]{.72\linewidth}\raggedright#1\end{varwidth}}
% Match the reference practice graph treatments; the toggle also controls ink.
\colorlet{keyMassColor}{blue!75!black}
\colorlet{keyCDFColor}{cyan!65!blue!80!black}
\colorlet{keyMassFill}{GTenHLProbability!32!white}
\colorlet{keyGuideColor}{keyMassColor!35}
\ifGTenHighlights\else
  \colorlet{keyMassColor}{black}
  \colorlet{keyCDFColor}{black}
  \colorlet{keyMassFill}{black!15}
  \colorlet{keyGuideColor}{gray}
\fi
\linespread{0.98}
\setlength{\parskip}{2pt}
\renewcommand{\arraystretch}{1.0}
\setlength{\emergencystretch}{2em}
\pgfplotsset{keyaxis/.style={axis lines=left,scaled ticks=false,
 xlabel={$\GTenHLValue{x}$ (\GTenHLValue{USD})},
 xticklabel={ $\GTenHLValue{\pgfmathprintnumber{\tick}}$},grid=major,grid style={gray!15},
 tick label style={font=\small},label style={font=\small},clip=false}}
\pgfplotsset{keyCombined/.style={width=.84\linewidth,height=4cm,scale only axis,
every axis plot/.append style={forget plot},axis lines=left,x axis line style={->},ymin=-.035,ymax=1.13,
ytick={0,.25,.5,.75,1},grid=major,
grid style={gray!15},tick label style={font=\small},
label style={font=\small},clip=false,
legend style={at={(.5,-.34)},anchor=north,legend columns=2,
font=\footnotesize,draw=none,cells={anchor=west}},
},keyMass/.style={keyMassColor,thick},
keyPhantom/.style={keyGuideColor,densely dashed,line width=.9pt},
keyStep/.style={keyCDFColor,thick},
keyIncluded/.style={only marks,mark=*,keyCDFColor,mark size=1.1pt},
keyExcluded/.style={only marks,mark=*,keyCDFColor,
mark options={fill=white},mark size=1.1pt}}
% Arguments: supported value, exact mass, cumulative height immediately before it.
% Only a vertical translation is applied: upper height = lower height + mass.
% Coordinates are saved inside their own axes, then joined after both axes end.
\usetikzlibrary{backgrounds,calc}
\tikzset{keyCorrespondence/.style={gray!70,line width=.25pt,-{Stealth[length=3.5pt,width=3pt]}}}
\newcommand{\keyCorrespondence}[1]{%
\begingroup
% Maximum sideways bow in physical units, independent of the axis scales.
\def\keyCorrespondenceBow{3mm}%
% Quadratic Bezier with midpoint control shifted by twice the bow;
% its equivalent cubic controls lie at 1/3 and 2/3, shifted by 4/3 of it.
\pgfmathsetlengthmacro{\keyControlOffset}{4/3*\keyCorrespondenceBow}%
% Paint behind the stems and markers, even though both axes are already defined.
\begin{pgfonlayer}{background}
\draw[keyCorrespondence] (keyMassMid-#1)
  .. controls ([xshift=\keyControlOffset]$(keyMassMid-#1)!1/3!(keyShadowMid-#1)$)
          and ([xshift=\keyControlOffset]$(keyMassMid-#1)!2/3!(keyShadowMid-#1)$)
  .. (keyShadowMid-#1);
\end{pgfonlayer}%
\endgroup}
\newcommand{\keyPhantom}[3]{%
\coordinate (keyShadowMid-#1) at (axis cs:#1,{#3+(#2)/2});
\addplot[keyPhantom] coordinates {(#1,#3)(#1,{#3+#2})};
\addplot[only marks,mark=*,keyGuideColor,mark size=1.8pt]
coordinates {(#1,{#3+#2})};}
\newcommand{\keyMass}[2]{%
\coordinate (keyMassMid-#1) at (axis cs:#1,{(#2)/2});
\addplot[keyMass] coordinates {(#1,0)(#1,#2)};
\addplot[only marks,mark=*,keyMassColor,mark size=1.5pt]
coordinates {(#1,#2)};}

\newcommand{\keyConnectionReading}{%
\textbf{Reading the comparison.} Both plotting areas use identical horizontal
limits, support ticks, numerical spacing, widths and vertical scales.
Each pale dashed shadow translates its lower \GTenHLProbability{PMF} stem upward
by $\HLCumulativeProbability{F_X(}\GTenHLValue{a^-}\HLCumulativeProbability{)}$ without changing its length:
$\HLCumulativeProbability{F_X(}\GTenHLValue{a}\HLCumulativeProbability{)}-\HLCumulativeProbability{F_X(}\GTenHLValue{a^-}\HLCumulativeProbability{)}=\GTenHLProbability{p_X(}\GTenHLValue{a}\GTenHLProbability{)}$.
Thin curved arrows join the stem and shadow midpoints, including the first
support value. They indicate correspondence. Stems and shadows show isolated
masses; the filled \HLCumulativeProbability{CDF} endpoint gives its value and
the open endpoint excludes the lower height.}

\title{Global Economics Exam Solutions\\Grade 10 --- Term 2 --- C1 and C2}
\author{}
\date{}
\makeatletter
\renewcommand{\@maketitle}{\begin{center}{\Large\bfseries\@title\par}\end{center}}
\renewcommand\subsection{\@startsection{subsection}{2}{\z@}{-1.4ex plus -.3ex}{.7ex}{\normalfont\large\bfseries}}
\makeatother
\hypersetup{pdftitle={Grade 10 Term 2 C1 and C2 Exam Solutions},pdfauthor={Global Economics}}
\begin{document}
\maketitle
\small
\setlength{\topsep}{3pt}
\setcounter{tocdepth}{1}
\phantomsection
\label{toc}
\tableofcontents
\section*{Learning objective}
Model discrete monetary outcomes in economic and financial contexts; derive and verify probability mass functions (PMFs), and construct and interpret cumulative distribution functions (CDFs) with tables and graphs.
\section*{Criteria assessment}
This key reproduces each alternate question before its worked solution. Questions 1--2 assess C1; Questions 3--6 assess C2. Show reasoning and exact probabilities. All contextual rules are stated in ordinary language.

This exam assesses C1 and C2. A \textbf{minimum of 4 out of 6 points for C1 and 4 out of 14 points for C2} is required to pass. The overall available total is 20 points. Show all relevant calculations and justify your answers. Partial credit may be awarded for correct procedures despite arithmetic errors; answers without demonstrated procedures receive no credit. A 1-point penalty per criterion applies if the student's name is missing or the assessment is submitted late.

\textbf{Independent passing thresholds: C1 requires at least 4 out of 6 points; C2 requires at least 4 out of 14 points. Points cannot transfer between criteria.} The overall available total is 20 points. An overall total alone does not establish a pass in both criteria. These are the passing thresholds for this assessment; no cumulative policy is specified.
\section*{Evaluated criteria}
\begin{description}
\item[C1] Characterises discrete random variables in economic and financial contexts by identifying their support, probability mass function and possible outcomes.
\item[C2] Interprets PMFs and CDFs by verifying their defining properties and using them to calculate point, interval, cumulative, and complementary probabilities in financial contexts.
\end{description}

\begin{center}
\begin{tabular}{|c|c|c|}
	\hline\textbf{\textnumero} & \textbf{C1} & \textbf{C2} \\ \hline
	\ref*{sec:q1} & 3 & --- \\\hline
	\ref*{sec:q2} & 3 & --- \\\hline
	\ref*{sec:q3} & --- & 2 \\\hline
	\ref*{sec:q4} & --- & 2 \\\hline
	\ref*{sec:q5} & --- & 4 \\\hline
	\ref*{sec:q6} & --- & 6 \\\hline
	\textbf{Total} & \textbf{6} & \textbf{14} \\\hline
\end{tabular}
\end{center}

\textbf{Marking policy.} Question and criterion totals are fixed by the brief. Task totals match the student exam; detailed partial-credit breakdowns below are proposed. Award stated fractions for demonstrated components, including valid alternative methods. Allow consistent follow-through after one arithmetic error when subsequent reasoning is correct. In C1 Task 2 (Questions 1--2) and C2 Task 1 (Questions 3--6), positivity (nonnegativity) and normalization checks carry no separate points; omitting these checks does not reduce the task score. Solving the normalization equation to derive probabilities in Question 5 remains assessed, as do its relationship checks. CDF verification retains its stated credit and must be demonstrated. Cap each task, question and criterion at its allocation.


\clearpage
\section*{Highlight legend}
\phantomsection\label{tab:highlight-legend}
\begin{center}
\begin{tabular}{@{}lll@{}}\toprule
\textbf{Symbol or notation} & \textbf{Concept} & \textbf{Brief meaning}\\\midrule
$\GTenHLOutcome{\omega}$ & \GTenHLOutcome{Outcome} & One complete result or category.\\
$\GTenHLValue{X},\GTenHLValue{x},\GTenHLValue{S_X}$ & \GTenHLValue{Variable and support} & Values and cutoffs in \GTenHLValue{USD}.\\
$\GTenHLSecond{N},\GTenHLSecond{Y}$ & \GTenHLSecond{Other quantities} & Counts; adjusted payout in USD.\\
$\GTenHLProbability{P},\GTenHLProbability{p_X(}\GTenHLValue{x}\GTenHLProbability{)}$ & \GTenHLProbability{Probability / PMF} & Unitless likelihood or point mass.\\
$\HLCumulativeProbability{F_X(}\GTenHLValue{x}\HLCumulativeProbability{)}$ & \HLCumulativeProbability{Cumulative probability / CDF} & Mass at or below a cutoff.\\
$\GTenHLParameter{a}$ & \GTenHLParameter{Unknown parameter} & Baseline mass to determine in Question 5.\\
\bottomrule\end{tabular}
\end{center}

\sectionScored{Bicycle tour contribution}{C1-\sml{3}}{sec:q1}

A cycling guide handles three bicycle tour bookings. The three tour bookings are independent. Each is either kept or cancelled, two mutually exclusive and exhaustive outcomes. The keep probabilities for tour bookings 1, 2 and 3 are 0.25, 0.75, 0.5, respectively. Each kept booking contributes 10 USD and each cancellation contributes no money. A fixed setup cost of 4 USD is then deducted once from the combined contribution. There are no other receipts or costs. Record the net contribution for the three tour bookings.

\Needspace{8\baselineskip}
\subsection*{Task 1 (a) --- question and worked solution (proposed allocation) \score{1}}
\label{sol:q1t1}

\noindent\textbf{\boldmath Define the random variable and its units, and identify its support.}\par\nopagebreak

Let $\GTenHLValue{X}$ be the net contribution from the three tour bookings, in \GTenHLValue{USD}. If $\GTenHLSecond{N}$ is the number kept,
\[\GTenHLValue{X}=\GTenHLValue{10}\GTenHLSecond{N}-\GTenHLValue{4}.\]
Since $N\in\{0,1,2,3\}$, the monetary values are $10(0)-4=-4,\quad 10(1)-4=6,\quad 10(2)-4=16,\quad 10(3)-4=26$ USD.

$\GTenHLValue{S_X}=\{\GTenHLValue{-4},\GTenHLValue{6},\GTenHLValue{16},\GTenHLValue{26}\}$ USD.
A negative contribution means the setup cost exceeds receipts.


\par\smallskip
\noindent\begin{minipage}{\linewidth}\centering
\textbf{Proposed partial credit.}\par\smallskip
\begin{tabular}{@{}l r@{}}\toprule
\CreditDescription{\textbf{Assessed step or criterion}} & \textbf{Points}\\\midrule
\CreditDescription{Contextual definition of $X$, monetary rule and units} & 0.5\\
\CreditDescription{Correct, complete support} & 0.5\\
\midrule
\CreditDescription{\textbf{Total}} & \textbf{1}\\
\bottomrule\end{tabular}
\end{minipage}\par\smallskip

\par\vspace{3mm}\secInicio[][sec:q1]

\clearpage
\subsection*{Task 2 (b) --- question and worked solution (proposed allocation) \score{1}}
\label{sol:q1t2}

\noindent\textbf{\boldmath List every complete outcome, preserving order where relevant; map each outcome to its probability and monetary value; and construct the grouped PMF.}\par\nopagebreak

Write $\GTenHLOutcome{K}$ for kept and $\GTenHLOutcome{C}$ for cancelled. Each triple follows the stated order. Before multiplying, find each complementary probability:
\[\begin{aligned}P(C_1)&=1-P(K_1)=1-0.25=\GTenHLProbability{0.75}\\
P(C_2)&=1-P(K_2)=1-0.75=\GTenHLProbability{0.25}\\
P(C_3)&=1-P(K_3)=1-0.5=\GTenHLProbability{0.5}.\end{aligned}\]
There are $2^3=8$ ordered complete outcomes. Independence permits products of three probabilities, each using its own trial probability.

\begin{center}\begin{tabular}{cccc}\toprule
Ordered outcome & Probability calculation & Exact probability & $\GTenHLValue{X}$ (USD)\\\midrule
$\GTenHLOutcome{KKK}$ & $\GTenHLProbability{0.25}\cdot\GTenHLProbability{0.75}\cdot\GTenHLProbability{0.5}$ & $\GTenHLProbability{3/32}=\GTenHLProbability{0.09375}=\GTenHLProbability{9.375\%}$ & $\GTenHLValue{26}$\\[4pt]
$\GTenHLOutcome{KKC}$ & $\GTenHLProbability{0.25}\cdot\GTenHLProbability{0.75}\cdot\GTenHLProbability{0.5}$ & $\GTenHLProbability{3/32}=\GTenHLProbability{0.09375}=\GTenHLProbability{9.375\%}$ & $\GTenHLValue{16}$\\[4pt]
$\GTenHLOutcome{KCK}$ & $\GTenHLProbability{0.25}\cdot\GTenHLProbability{0.25}\cdot\GTenHLProbability{0.5}$ & $\GTenHLProbability{1/32}=\GTenHLProbability{0.03125}=\GTenHLProbability{3.125\%}$ & $\GTenHLValue{16}$\\[4pt]
$\GTenHLOutcome{KCC}$ & $\GTenHLProbability{0.25}\cdot\GTenHLProbability{0.25}\cdot\GTenHLProbability{0.5}$ & $\GTenHLProbability{1/32}=\GTenHLProbability{0.03125}=\GTenHLProbability{3.125\%}$ & $\GTenHLValue{6}$\\[4pt]
$\GTenHLOutcome{CKK}$ & $\GTenHLProbability{0.75}\cdot\GTenHLProbability{0.75}\cdot\GTenHLProbability{0.5}$ & $\GTenHLProbability{9/32}=\GTenHLProbability{0.28125}=\GTenHLProbability{28.125\%}$ & $\GTenHLValue{16}$\\[4pt]
$\GTenHLOutcome{CKC}$ & $\GTenHLProbability{0.75}\cdot\GTenHLProbability{0.75}\cdot\GTenHLProbability{0.5}$ & $\GTenHLProbability{9/32}=\GTenHLProbability{0.28125}=\GTenHLProbability{28.125\%}$ & $\GTenHLValue{6}$\\[4pt]
$\GTenHLOutcome{CCK}$ & $\GTenHLProbability{0.75}\cdot\GTenHLProbability{0.25}\cdot\GTenHLProbability{0.5}$ & $\GTenHLProbability{3/32}=\GTenHLProbability{0.09375}=\GTenHLProbability{9.375\%}$ & $\GTenHLValue{6}$\\[4pt]
$\GTenHLOutcome{CCC}$ & $\GTenHLProbability{0.75}\cdot\GTenHLProbability{0.25}\cdot\GTenHLProbability{0.5}$ & $\GTenHLProbability{3/32}=\GTenHLProbability{0.09375}=\GTenHLProbability{9.375\%}$ & $\GTenHLValue{-4}$\\[4pt]
\bottomrule\end{tabular}\end{center}
Exhaustiveness check: $(0.25+0.75)(0.75+0.25)(0.5+0.5)=1$; the complete outcome probabilities sum to 1. Outcomes are not equally likely.

\par\vspace{3mm}\secInicio[][sec:q1]
\clearpage\textbf{Question 1 --- Task 2 (b) continued: grouped \GTenHLProbability{PMF}}

For one complete outcome $\omega$, $P(\{\omega\})$ is its probability. Add the probabilities of all outcomes mapping to the same monetary value: $p_X(x)=\sum_{\omega:X(\omega)=x}P(\{\omega\})$.

\begin{center}\begin{tabular}{@{}p{.30\linewidth}c p{.54\linewidth}@{}}\toprule
Complete outcomes & $\GTenHLValue{x}$ (USD) & Grouped \GTenHLProbability{PMF}\\\midrule
$\GTenHLOutcome{CCC}$ & $\GTenHLValue{-4}$ & $\begin{aligned}[t]p_X(\GTenHLValue{-4})&=\GTenHLProbability{3/32}\\&=\GTenHLProbability{3/32}=\GTenHLProbability{0.09375}=\GTenHLProbability{9.375\%}\end{aligned}$\\[4pt]
$\GTenHLOutcome{KCC}$, $\GTenHLOutcome{CKC}$, $\GTenHLOutcome{CCK}$ & $\GTenHLValue{6}$ & $\begin{aligned}[t]p_X(\GTenHLValue{6})&=\GTenHLProbability{1/32}+\GTenHLProbability{9/32}+\GTenHLProbability{3/32}\\&=\GTenHLProbability{13/32}=\GTenHLProbability{0.40625}=\GTenHLProbability{40.625\%}\end{aligned}$\\[4pt]
$\GTenHLOutcome{KKC}$, $\GTenHLOutcome{KCK}$, $\GTenHLOutcome{CKK}$ & $\GTenHLValue{16}$ & $\begin{aligned}[t]p_X(\GTenHLValue{16})&=\GTenHLProbability{3/32}+\GTenHLProbability{1/32}+\GTenHLProbability{9/32}\\&=\GTenHLProbability{13/32}=\GTenHLProbability{0.40625}=\GTenHLProbability{40.625\%}\end{aligned}$\\[4pt]
$\GTenHLOutcome{KKK}$ & $\GTenHLValue{26}$ & $\begin{aligned}[t]p_X(\GTenHLValue{26})&=\GTenHLProbability{3/32}\\&=\GTenHLProbability{3/32}=\GTenHLProbability{0.09375}=\GTenHLProbability{9.375\%}\end{aligned}$\\[4pt]
\bottomrule\end{tabular}\end{center}
\Needspace{5\baselineskip}\textbf{PMF checks.} All masses are positive. Normalization: $(3/32+13/32+13/32+3/32)=1$. Set $p_X(x)=0$ outside the support; gaps carry no mass.

\textbf{Scoring note.} These explanatory positivity and normalization checks carry no separate points; omitting them does not reduce the task score.


\par\smallskip
\noindent\begin{minipage}{\linewidth}\centering
\textbf{Proposed partial credit.}\par\smallskip
\begin{tabular}{@{}l r@{}}\toprule
\CreditDescription{\textbf{Assessed step or criterion}} & \textbf{Points}\\\midrule
\CreditDescription{Complete ordered outcomes} & 0.25\\
\CreditDescription{Correct outcome probability calculations} & 0.25\\
\CreditDescription{Correct monetary mapping for every outcome} & 0.25\\
\CreditDescription{Correct grouping and grouped PMF masses} & 0.25\\
\midrule
\CreditDescription{\textbf{Total}} & \textbf{1}\\
\bottomrule\end{tabular}
\end{minipage}\par\smallskip

\par\vspace{3mm}\secInicio[][sec:q1]

\Needspace{8\baselineskip}
\subsection*{Task 3 (c) --- question and worked solution (proposed allocation) \score{1}}
\label{sol:q1t3}

\noindent\textbf{\boldmath Calculate and interpret the probability that the net contribution is exactly 16 USD.}\par\nopagebreak

$P(\GTenHLValue{X}=\GTenHLValue{16})=p_X(\GTenHLValue{16})=\GTenHLProbability{13/32}=\GTenHLProbability{13/32}=\GTenHLProbability{0.40625}=\GTenHLProbability{40.625\%}$.

The net contribution is exactly 16 USD with probability $\GTenHLProbability{13/32}$. This occurs exactly for $\GTenHLOutcome{KKC}$, $\GTenHLOutcome{KCK}$, $\GTenHLOutcome{CKK}$.


\par\smallskip
\noindent\begin{minipage}{\linewidth}\centering
\textbf{Proposed partial credit.}\par\smallskip
\begin{tabular}{@{}l r@{}}\toprule
\CreditDescription{\textbf{Assessed step or criterion}} & \textbf{Points}\\\midrule
\CreditDescription{Selecting the correct mass at $X=16$} & 0.5\\
\CreditDescription{The contextual interpretation} & 0.5\\
\midrule
\CreditDescription{\textbf{Total}} & \textbf{1}\\
\bottomrule\end{tabular}
\end{minipage}\par\smallskip

Question total: 3 points for C1.

\par\vspace{3mm}\secInicio[][sec:q1]

\clearpage
\sectionScored{Two packaging return rewards}{C1-\sml{3}}{sec:q2}

A grocery shop receives one return reward from each of two packaging suppliers. The packaging suppliers award payments independently. Each has exactly three mutually exclusive possible outcomes: no payment, a 5 USD payment, or an 12 USD payment. For the first provider these probabilities are 0.15, 0.7, 0.15, respectively. For the second provider they are 0.15, 0.3, 0.55, respectively. Add both payments; there are no other adjustments. Record the total payment, keeping first-provider then second-provider order.

\Needspace{8\baselineskip}
\subsection*{Task 1 (a) --- question and worked solution (proposed allocation) \score{1}}
\label{sol:q2t1}

\noindent\textbf{\boldmath Define the random variable and its units, and identify its support.}\par\nopagebreak

Let $\GTenHLValue{X}$ be the sum of the two providers' payments, in \GTenHLValue{USD}. Each payment belongs to $\{0,5,12\}$ USD. Add the two payments for each ordered outcome:
\[\begin{array}{lll}0+0=0&0+5=5&0+12=12\\
5+0=5&5+5=10&5+12=17\\
12+0=12&12+5=17&12+12=24.\end{array}\]
Repeated sums appear only once in the support. $\GTenHLValue{S_X}=\{\GTenHLValue{0},\GTenHLValue{5},\GTenHLValue{10},\GTenHLValue{12},\GTenHLValue{17},\GTenHLValue{24}\}$ USD.


\par\smallskip
\noindent\begin{minipage}{\linewidth}\centering
\textbf{Proposed partial credit.}\par\smallskip
\begin{tabular}{@{}l r@{}}\toprule
\CreditDescription{\textbf{Assessed step or criterion}} & \textbf{Points}\\\midrule
\CreditDescription{Contextual definition of $X$, monetary rule and units} & 0.5\\
\CreditDescription{Correct, complete support} & 0.5\\
\midrule
\CreditDescription{\textbf{Total}} & \textbf{1}\\
\bottomrule\end{tabular}
\end{minipage}\par\smallskip

\par\vspace{3mm}\secInicio[][sec:q2]

\clearpage
\subsection*{Task 2 (b) --- question and worked solution (proposed allocation) \score{1}}
\label{sol:q2t2}

\noindent\textbf{\boldmath List every complete outcome, preserving order where relevant; map each outcome to its probability and monetary value; and construct the grouped PMF.}\par\nopagebreak

Each provider has three possible outcomes, giving $3^2=9$ ordered pairs. Preserve first-provider then second-provider order. Independence permits multiplication of two marginal probabilities.

\begin{center}\begin{tabular}{cccc}\toprule
Ordered outcome & Probability calculation & Exact probability & $\GTenHLValue{X}$ (USD)\\\midrule
$\GTenHLOutcome{(0,0)}$ & $\GTenHLProbability{0.15}\cdot\GTenHLProbability{0.15}$ & $\GTenHLProbability{9/400}=\GTenHLProbability{0.0225}=\GTenHLProbability{2.25\%}$ & $\GTenHLValue{0}$\\[4pt]
$\GTenHLOutcome{(0,5)}$ & $\GTenHLProbability{0.15}\cdot\GTenHLProbability{0.3}$ & $\GTenHLProbability{9/200}=\GTenHLProbability{0.045}=\GTenHLProbability{4.5\%}$ & $\GTenHLValue{5}$\\[4pt]
$\GTenHLOutcome{(0,12)}$ & $\GTenHLProbability{0.15}\cdot\GTenHLProbability{0.55}$ & $\GTenHLProbability{33/400}=\GTenHLProbability{0.0825}=\GTenHLProbability{8.25\%}$ & $\GTenHLValue{12}$\\[4pt]
$\GTenHLOutcome{(5,0)}$ & $\GTenHLProbability{0.7}\cdot\GTenHLProbability{0.15}$ & $\GTenHLProbability{21/200}=\GTenHLProbability{0.105}=\GTenHLProbability{10.5\%}$ & $\GTenHLValue{5}$\\[4pt]
$\GTenHLOutcome{(5,5)}$ & $\GTenHLProbability{0.7}\cdot\GTenHLProbability{0.3}$ & $\GTenHLProbability{21/100}=\GTenHLProbability{0.21}=\GTenHLProbability{21\%}$ & $\GTenHLValue{10}$\\[4pt]
$\GTenHLOutcome{(5,12)}$ & $\GTenHLProbability{0.7}\cdot\GTenHLProbability{0.55}$ & $\GTenHLProbability{77/200}=\GTenHLProbability{0.385}=\GTenHLProbability{38.5\%}$ & $\GTenHLValue{17}$\\[4pt]
$\GTenHLOutcome{(12,0)}$ & $\GTenHLProbability{0.15}\cdot\GTenHLProbability{0.15}$ & $\GTenHLProbability{9/400}=\GTenHLProbability{0.0225}=\GTenHLProbability{2.25\%}$ & $\GTenHLValue{12}$\\[4pt]
$\GTenHLOutcome{(12,5)}$ & $\GTenHLProbability{0.15}\cdot\GTenHLProbability{0.3}$ & $\GTenHLProbability{9/200}=\GTenHLProbability{0.045}=\GTenHLProbability{4.5\%}$ & $\GTenHLValue{17}$\\[4pt]
$\GTenHLOutcome{(12,12)}$ & $\GTenHLProbability{0.15}\cdot\GTenHLProbability{0.55}$ & $\GTenHLProbability{33/400}=\GTenHLProbability{0.0825}=\GTenHLProbability{8.25\%}$ & $\GTenHLValue{24}$\\[4pt]
\bottomrule\end{tabular}\end{center}
Exhaustiveness check: $(0.15+0.7+0.15)(0.15+0.3+0.55)=1$; the complete outcome probabilities sum to 1. Outcomes are not equally likely.

\par\vspace{3mm}\secInicio[][sec:q2]
\clearpage\textbf{Question 2 --- Task 2 (b) continued: grouped \GTenHLProbability{PMF}}

For one complete outcome $\omega$, $P(\{\omega\})$ is its probability. Add the probabilities of all outcomes mapping to the same monetary value: $p_X(x)=\sum_{\omega:X(\omega)=x}P(\{\omega\})$.

\begin{center}\begin{tabular}{@{}p{.30\linewidth}c p{.54\linewidth}@{}}\toprule
Complete outcomes & $\GTenHLValue{x}$ (USD) & Grouped \GTenHLProbability{PMF}\\\midrule
$\GTenHLOutcome{(0,0)}$ & $\GTenHLValue{0}$ & $\begin{aligned}[t]p_X(\GTenHLValue{0})&=\GTenHLProbability{9/400}\\&=\GTenHLProbability{9/400}=\GTenHLProbability{0.0225}=\GTenHLProbability{2.25\%}\end{aligned}$\\[4pt]
$\GTenHLOutcome{(0,5)}$, $\GTenHLOutcome{(5,0)}$ & $\GTenHLValue{5}$ & $\begin{aligned}[t]p_X(\GTenHLValue{5})&=\GTenHLProbability{9/200}+\GTenHLProbability{21/200}\\&=\GTenHLProbability{3/20}=\GTenHLProbability{0.15}=\GTenHLProbability{15\%}\end{aligned}$\\[4pt]
$\GTenHLOutcome{(5,5)}$ & $\GTenHLValue{10}$ & $\begin{aligned}[t]p_X(\GTenHLValue{10})&=\GTenHLProbability{21/100}\\&=\GTenHLProbability{21/100}=\GTenHLProbability{0.21}=\GTenHLProbability{21\%}\end{aligned}$\\[4pt]
$\GTenHLOutcome{(0,12)}$, $\GTenHLOutcome{(12,0)}$ & $\GTenHLValue{12}$ & $\begin{aligned}[t]p_X(\GTenHLValue{12})&=\GTenHLProbability{33/400}+\GTenHLProbability{9/400}\\&=\GTenHLProbability{21/200}=\GTenHLProbability{0.105}=\GTenHLProbability{10.5\%}\end{aligned}$\\[4pt]
$\GTenHLOutcome{(5,12)}$, $\GTenHLOutcome{(12,5)}$ & $\GTenHLValue{17}$ & $\begin{aligned}[t]p_X(\GTenHLValue{17})&=\GTenHLProbability{77/200}+\GTenHLProbability{9/200}\\&=\GTenHLProbability{43/100}=\GTenHLProbability{0.43}=\GTenHLProbability{43\%}\end{aligned}$\\[4pt]
$\GTenHLOutcome{(12,12)}$ & $\GTenHLValue{24}$ & $\begin{aligned}[t]p_X(\GTenHLValue{24})&=\GTenHLProbability{33/400}\\&=\GTenHLProbability{33/400}=\GTenHLProbability{0.0825}=\GTenHLProbability{8.25\%}\end{aligned}$\\[4pt]
\bottomrule\end{tabular}\end{center}
\Needspace{5\baselineskip}\textbf{PMF checks.} All masses are positive. Normalization: $(9/400+3/20+21/100+21/200+43/100+33/400)=1$. Set $p_X(x)=0$ outside the support; gaps carry no mass.

\textbf{Scoring note.} These explanatory positivity and normalization checks carry no separate points; omitting them does not reduce the task score.


\par\smallskip
\noindent\begin{minipage}{\linewidth}\centering
\textbf{Proposed partial credit.}\par\smallskip
\begin{tabular}{@{}l r@{}}\toprule
\CreditDescription{\textbf{Assessed step or criterion}} & \textbf{Points}\\\midrule
\CreditDescription{Complete ordered outcomes} & 0.25\\
\CreditDescription{Correct outcome probability calculations} & 0.25\\
\CreditDescription{Correct monetary mapping for every outcome} & 0.25\\
\CreditDescription{Correct grouping and grouped PMF masses} & 0.25\\
\midrule
\CreditDescription{\textbf{Total}} & \textbf{1}\\
\bottomrule\end{tabular}
\end{minipage}\par\smallskip

\par\vspace{3mm}\secInicio[][sec:q2]

\Needspace{8\baselineskip}
\subsection*{Task 3 (c) --- question and worked solution (proposed allocation) \score{1}}
\label{sol:q2t3}

\noindent\textbf{\boldmath Calculate and interpret the probability that the total payment is exactly 17 USD.}\par\nopagebreak

$P(\GTenHLValue{X}=\GTenHLValue{17})=p_X(\GTenHLValue{17})=\GTenHLProbability{77/200}+\GTenHLProbability{9/200}=\GTenHLProbability{43/100}=\GTenHLProbability{0.43}=\GTenHLProbability{43\%}$.

The total payment is exactly 17 USD with probability $\GTenHLProbability{43/100}$. This occurs exactly for $\GTenHLOutcome{(5,12)}$, $\GTenHLOutcome{(12,5)}$.


\par\smallskip
\noindent\begin{minipage}{\linewidth}\centering
\textbf{Proposed partial credit.}\par\smallskip
\begin{tabular}{@{}l r@{}}\toprule
\CreditDescription{\textbf{Assessed step or criterion}} & \textbf{Points}\\\midrule
\CreditDescription{Selecting the correct mass at $X=17$} & 0.5\\
\CreditDescription{The contextual interpretation} & 0.5\\
\midrule
\CreditDescription{\textbf{Total}} & \textbf{1}\\
\bottomrule\end{tabular}
\end{minipage}\par\smallskip

Question total: 3 points for C1. Criterion C1 total: 6 points; independent pass: at least 4 points.

\par\vspace{3mm}\secInicio[][sec:q2]

\clearpage
\sectionScored{Festival cloakroom charges}{C2-\sml{2}}{sec:q3}

A festival accepts one cloakroom item under exactly one of eight exhaustive, mutually exclusive attendee categories: Volunteers, Performers, Students, Seniors, Spectators, Vendors, Sponsors or Guests. Their probabilities are 0.125, 0.125, 0.125, 0.125, 0.125, 0.125, 0.125, 0.125, respectively. The corresponding fixed charges are 0, 0, 5, 5, 12, 12, 24, 24 USD, respectively. No other charges apply. Record the charge for this transaction.

\Needspace{8\baselineskip}
\subsection*{Task 1 --- question and worked solution (proposed allocation) \score{0.5}}
\label{sol:q3t1}

\noindent\textbf{\boldmath Define the random variable and support; derive the grouped PMF; produce its table and a labeled graph with actual monetary spacing.}\par\nopagebreak

Let $\GTenHLValue{X}$ be the charge for one transaction, in \GTenHLValue{USD}. $\GTenHLValue{S_X}=\{\GTenHLValue{0},\GTenHLValue{5},\GTenHLValue{12},\GTenHLValue{24}\}$ USD.

The eight distinct categories map to four monetary values; probabilities are supplied, not assumed equal.

\begin{center}\begin{tabular}{ccc}\toprule
Category & Given probability & $\GTenHLValue{X}$ (USD)\\\midrule
\GTenHLOutcome{Volunteers} & $\GTenHLProbability{0.125}=\GTenHLProbability{1/8}$ & $\GTenHLValue{0}$\\[4pt]
\GTenHLOutcome{Performers} & $\GTenHLProbability{0.125}=\GTenHLProbability{1/8}$ & $\GTenHLValue{0}$\\[4pt]
\GTenHLOutcome{Students} & $\GTenHLProbability{0.125}=\GTenHLProbability{1/8}$ & $\GTenHLValue{5}$\\[4pt]
\GTenHLOutcome{Seniors} & $\GTenHLProbability{0.125}=\GTenHLProbability{1/8}$ & $\GTenHLValue{5}$\\[4pt]
\GTenHLOutcome{Spectators} & $\GTenHLProbability{0.125}=\GTenHLProbability{1/8}$ & $\GTenHLValue{12}$\\[4pt]
\GTenHLOutcome{Vendors} & $\GTenHLProbability{0.125}=\GTenHLProbability{1/8}$ & $\GTenHLValue{12}$\\[4pt]
\GTenHLOutcome{Sponsors} & $\GTenHLProbability{0.125}=\GTenHLProbability{1/8}$ & $\GTenHLValue{24}$\\[4pt]
\GTenHLOutcome{Guests} & $\GTenHLProbability{0.125}=\GTenHLProbability{1/8}$ & $\GTenHLValue{24}$\\[4pt]
\bottomrule\end{tabular}\end{center}

\par\vspace{3mm}\secInicio[][sec:q3]
\clearpage\textbf{Question 3 --- Task 1 continued: grouped \GTenHLProbability{PMF} and graph}


\begin{center}\begin{tabular}{@{}p{.30\linewidth}c p{.54\linewidth}@{}}\toprule
Complete outcomes & $\GTenHLValue{x}$ (USD) & Grouped \GTenHLProbability{PMF}\\\midrule
\GTenHLOutcome{Volunteers}, \GTenHLOutcome{Performers} & $\GTenHLValue{0}$ & $\begin{aligned}[t]p_X(\GTenHLValue{0})&=\GTenHLProbability{1/8}+\GTenHLProbability{1/8}\\&=\GTenHLProbability{1/4}=\GTenHLProbability{0.25}=\GTenHLProbability{25\%}\end{aligned}$\\[4pt]
\GTenHLOutcome{Students}, \GTenHLOutcome{Seniors} & $\GTenHLValue{5}$ & $\begin{aligned}[t]p_X(\GTenHLValue{5})&=\GTenHLProbability{1/8}+\GTenHLProbability{1/8}\\&=\GTenHLProbability{1/4}=\GTenHLProbability{0.25}=\GTenHLProbability{25\%}\end{aligned}$\\[4pt]
\GTenHLOutcome{Spectators}, \GTenHLOutcome{Vendors} & $\GTenHLValue{12}$ & $\begin{aligned}[t]p_X(\GTenHLValue{12})&=\GTenHLProbability{1/8}+\GTenHLProbability{1/8}\\&=\GTenHLProbability{1/4}=\GTenHLProbability{0.25}=\GTenHLProbability{25\%}\end{aligned}$\\[4pt]
\GTenHLOutcome{Sponsors}, \GTenHLOutcome{Guests} & $\GTenHLValue{24}$ & $\begin{aligned}[t]p_X(\GTenHLValue{24})&=\GTenHLProbability{1/8}+\GTenHLProbability{1/8}\\&=\GTenHLProbability{1/4}=\GTenHLProbability{0.25}=\GTenHLProbability{25\%}\end{aligned}$\\[4pt]
\bottomrule\end{tabular}\end{center}
\begin{center}\begin{tikzpicture}\begin{axis}[keyaxis,width=.86\linewidth,height=5cm,xmin=-2,xmax=26,xtick={0,5,12,24},xticklabel style={rotate=45,anchor=north east},ymin=0,ymax=.7,ylabel={$\GTenHLProbability{p_X(x)}$}]
\draw[fill=keyMassFill,draw=keyMassColor] (axis cs:-0.225,0) rectangle (axis cs:0.225,1/4);
\draw[fill=keyMassFill,draw=keyMassColor] (axis cs:4.775,0) rectangle (axis cs:5.225,1/4);
\draw[fill=keyMassFill,draw=keyMassColor] (axis cs:11.775,0) rectangle (axis cs:12.225,1/4);
\draw[fill=keyMassFill,draw=keyMassColor] (axis cs:23.775,0) rectangle (axis cs:24.225,1/4);
\end{axis}\end{tikzpicture}\end{center}
\Needspace{5\baselineskip}\textbf{PMF checks.} All masses are positive. Normalization: $(1/4+1/4+1/4+1/4)=1$. Set $p_X(x)=0$ outside the support; gaps carry no mass.

\textbf{Scoring note.} These explanatory positivity and normalization checks carry no separate points; omitting them does not reduce the task score.


\par\smallskip
\noindent\begin{minipage}{\linewidth}\centering
\textbf{Proposed partial credit.}\par\smallskip
\begin{tabular}{@{}l r@{}}\toprule
\CreditDescription{\textbf{Assessed step or criterion}} & \textbf{Points}\\\midrule
\CreditDescription{Variable, units and support} & 0.125\\
\CreditDescription{Category mapping, grouping and the PMF table} & 0.25\\
\CreditDescription{The labeled, correctly spaced graph} & 0.125\\
\midrule
\CreditDescription{\textbf{Total}} & \textbf{0.5}\\
\bottomrule\end{tabular}
\end{minipage}\par\smallskip

\par\vspace{3mm}\secInicio[][sec:q3]

\clearpage
\textbf{Question 3 --- continued}\quad\hyperref[sec:q3]{Original description}

\Needspace{8\baselineskip}
\subsection*{Task 2 --- question and worked solution (proposed allocation) \score{1}}
\label{sol:q3t2}

\noindent\textbf{\boldmath Construct the complete CDF table, including both tails, and draw a labeled step graph. Use open and filled endpoints correctly and preserve actual monetary spacing.}\par\nopagebreak

For $\HLCumulativeProbability{F_X(x)}=P(X\leq x)$, add the masses in increasing monetary order.


\begin{center}\begin{tabular}{@{}l l c@{}}\toprule
Cutoff interval (USD) & Cumulative calculation & $\HLCumulativeProbability{F_X(x)}$\\\midrule
$x<\GTenHLValue{0}$ & No included support values & $\HLCumulativeProbability{0}$\\[4pt]
$\GTenHLValue{0}\leq x<\GTenHLValue{5}$ & $\begin{aligned}[t]&\GTenHLProbability{1/4}\end{aligned}$ & $\HLCumulativeProbability{1/4}$\\[4pt]
$\GTenHLValue{5}\leq x<\GTenHLValue{12}$ & $\begin{aligned}[t]&\GTenHLProbability{1/4}+\GTenHLProbability{1/4}\end{aligned}$ & $\HLCumulativeProbability{1/2}$\\[4pt]
$\GTenHLValue{12}\leq x<\GTenHLValue{24}$ & $\begin{aligned}[t]&\GTenHLProbability{1/4}+\GTenHLProbability{1/4}+\GTenHLProbability{1/4}\end{aligned}$ & $\HLCumulativeProbability{3/4}$\\[4pt]
$\GTenHLValue{24}\leq x$ & $\begin{aligned}[t]&\GTenHLProbability{1/4}+\GTenHLProbability{1/4}+\GTenHLProbability{1/4}+\GTenHLProbability{1/4}\end{aligned}$ & $\HLCumulativeProbability{1}$\\[4pt]
\bottomrule\end{tabular}\end{center}
\begin{center}\begin{tikzpicture}\begin{axis}[keyaxis,width=.86\linewidth,height=7cm,xmin=-2,xmax=26,xtick={0,5,12,24},xticklabel style={rotate=45,anchor=north east},ymin=-.07,ymax=1.1,ytick={0,.25,.5,.75,1},ylabel={$\HLCumulativeProbability{F_X(x)}$}]
\addplot[keyCDFColor,thick,<-] coordinates {(-2,0)(0,0)};
\addplot[keyCDFColor,thick] coordinates {(0,1/4)(5,1/4)};
\addplot[keyCDFColor,thick] coordinates {(5,1/2)(12,1/2)};
\addplot[keyCDFColor,thick] coordinates {(12,3/4)(24,3/4)};
\addplot[keyCDFColor,thick,->] coordinates {(24,1)(26,1)};
\addplot[keyCDFColor,only marks,mark=*,mark size=1pt,mark options={fill=white}] coordinates {(0,0)(5,1/4)(12,1/2)(24,3/4)};
\addplot[keyCDFColor,only marks,mark=*,mark size=1pt] coordinates {(0,1/4)(5,1/2)(12,3/4)(24,1)};
\end{axis}\end{tikzpicture}\end{center}
\par\vspace{3mm}\secInicio[][sec:q3]
\clearpage\textbf{Question 3 --- Task 2 continued: comparison}

\textbf{Same distribution: CDF above, PMF below.}
\begin{center}\begin{minipage}{\linewidth}\centering\begin{tikzpicture}
\pgfplotsset{keySupport/.style={xmin=-2,xmax=26,xtick={0,5,12,24},xticklabel style={rotate=45,anchor=north east}}}
\begin{axis}[keyCombined,keySupport,name=keyTop,anchor=south west,ylabel={$\HLCumulativeProbability{F_X(x)}$}]
\keyPhantom{0}{1/4}{0}
\keyPhantom{5}{1/4}{1/4}
\keyPhantom{12}{1/4}{1/2}
\keyPhantom{24}{1/4}{3/4}
\addplot[keyStep,<-] coordinates {(-2,0)(0,0)};
\addplot[keyStep] coordinates {(0,1/4)(5,1/4)};
\addplot[keyStep] coordinates {(5,1/2)(12,1/2)};
\addplot[keyStep] coordinates {(12,3/4)(24,3/4)};
\addplot[keyStep,->] coordinates {(24,1)(26,1)};
\addplot[keyExcluded] coordinates {(0,0)(5,1/4)(12,1/2)(24,3/4)};
\addplot[keyIncluded] coordinates {(0,1/4)(5,1/2)(12,3/4)(24,1)};
\end{axis}
\begin{axis}[keyCombined,keySupport,at={(keyTop.south west)},anchor=north west,yshift=-1.35cm,ylabel={$\GTenHLProbability{p_X(x)}$}]
\keyMass{0}{1/4}
\keyMass{5}{1/4}
\keyMass{12}{1/4}
\keyMass{24}{1/4}
\addlegendimage{keyStep}\addlegendentry{Top: CDF steps}
\addlegendimage{keyPhantom,mark=*}\addlegendentry{Top: PMF shadows (jumps)}
\addlegendimage{keyMass,mark=*}\addlegendentry{Bottom: PMF stems}
\addlegendimage{keyCorrespondence}\addlegendentry{Midpoint correspondence}
\end{axis}
\keyCorrespondence{0}
\keyCorrespondence{5}
\keyCorrespondence{12}
\keyCorrespondence{24}
\end{tikzpicture}\end{minipage}\end{center}\keyConnectionReading

\textbf{CDF checks.}\begin{itemize}\item The intervals cover every real cutoff; all heights lie in $[0,1]$ and never decrease.\item Each segment includes its left endpoint and excludes its right endpoint. Filled endpoints give the values at jumps, so the CDF is right-continuous.\item The lower-tail limit is 0 and the upper-tail limit is 1; arrows extend both tails.\item Successive differences recover the PMF: $1/4-0=1/4,\quad 1/2-1/4=1/4,\quad 3/4-1/2=1/4,\quad 1-3/4=1/4$. Thus $F_X(a)-F_X(a^-)=p_X(a)$; unsupported cutoffs have zero jump and zero mass.\end{itemize}


\par\smallskip
\noindent\begin{minipage}{\linewidth}\centering
\textbf{Proposed partial credit.}\par\smallskip
\begin{tabular}{@{}l r@{}}\toprule
\CreditDescription{\textbf{Assessed step or criterion}} & \textbf{Points}\\\midrule
\CreditDescription{Cumulative sums} & 0.25\\
\CreditDescription{The complete interval table and tails} & 0.25\\
\CreditDescription{The labeled, correctly spaced step graph and endpoints} & 0.25\\
\CreditDescription{The CDF and jump checks} & 0.25\\
\midrule
\CreditDescription{\textbf{Total}} & \textbf{1}\\
\bottomrule\end{tabular}
\end{minipage}\par\smallskip

\par\vspace{3mm}\secInicio[][sec:q3]

\clearpage
\textbf{Question 3 --- continued}\quad\hyperref[sec:q3]{Original description}

\Needspace{8\baselineskip}
\subsection*{Task 3 (c) --- question and worked solution (proposed allocation) \score{0.5}}
\label{sol:q3t3}

\noindent\textbf{\boldmath Calculate and interpret (i) the probability that the charge is less than 12 USD and (ii) the probability that it is at least 24 USD.}\par\nopagebreak

\textbf{Allocation:} (i) 0.25 points; (ii) 0.25 points.

\textbf{(i)} Include only $\GTenHLValue{0},\GTenHLValue{5}$ USD; exclude the boundary mass at 12 USD.
\[\begin{aligned}P(\GTenHLValue{X}<\GTenHLValue{12})&=p_X(\GTenHLValue{0})+p_X(\GTenHLValue{5})\\&=\GTenHLProbability{1/4}+\GTenHLProbability{1/4}\\&=\GTenHLProbability{1/2}=\GTenHLProbability{0.5}=\GTenHLProbability{50\%}.\end{aligned}\]
The charge for one case is below 12 USD with probability $\GTenHLProbability{1/2}$, approximately 50.00\%.

\textbf{(ii)} Include only $\GTenHLValue{24}$ USD; include the boundary mass at 24 USD.
\[\begin{aligned}P(\GTenHLValue{X}\geq\GTenHLValue{24})&=p_X(\GTenHLValue{24})\\&=\GTenHLProbability{1/4}\\&=\GTenHLProbability{1/4}=\GTenHLProbability{0.25}=\GTenHLProbability{25\%}.\end{aligned}\]
The charge for one case is at least 24 USD with probability $\GTenHLProbability{1/4}$, approximately 25.00\%.



\par\smallskip
\noindent\begin{minipage}{\linewidth}\centering
\textbf{Proposed partial credit.}\par\smallskip
\begin{tabular}{@{}l r@{}}\toprule
\CreditDescription{\textbf{Assessed step or criterion}} & \textbf{Points}\\\midrule
\CreditDescription{(i): selecting exactly the included PMF masses with the correct boundary} & 0.125\\
\CreditDescription{(i): summing them correctly} & 0.0625\\
\CreditDescription{(i): the contextual interpretation} & 0.0625\\
\CreditDescription{(ii): selecting exactly the included PMF masses with the correct boundary} & 0.125\\
\CreditDescription{(ii): summing them correctly} & 0.0625\\
\CreditDescription{(ii): the contextual interpretation} & 0.0625\\
\midrule
\CreditDescription{\textbf{Total}} & \textbf{0.5}\\
\bottomrule\end{tabular}
\end{minipage}\par\smallskip

Each part totals 0.25 points.

\par\vspace{3mm}\secInicio[][sec:q3]

\textbf{Question total: $0.5+1+0.5=2$ points for C2.}

\clearpage
\sectionScored{Theatre audience reimbursements}{C2-\sml{2}}{sec:q4}

A theatre classifies one audience complaint in exactly one of eight exhaustive, mutually exclusive categories: No adjustment, Apology, Seat change, Programme, Partial refund, Replacement, Voucher or Cash. These eight named categories are the distinct outcomes in the sample space. Their probabilities are 0.07, 0.09, 0.13, 0.09, 0.18, 0.22, 0.08, 0.14, respectively; do not assume equal likelihood. The corresponding payouts are 0, 0, 7, 7, 17, 17, 26, 26 USD, respectively. Count the voucher at its face value as a payout. Each payout is fixed for its category, and no additional payments apply. Record the payout for one case.

\Needspace{8\baselineskip}
\subsection*{Task 1 --- question and worked solution (proposed allocation) \score{0.5}}
\label{sol:q4t1}

\noindent\textbf{\boldmath Define the random variable and support; derive the grouped PMF; produce its table and a labeled graph with actual monetary spacing.}\par\nopagebreak

Let $\GTenHLValue{X}$ be the payout for one case, in \GTenHLValue{USD}. $\GTenHLValue{S_X}=\{\GTenHLValue{0},\GTenHLValue{7},\GTenHLValue{17},\GTenHLValue{26}\}$ USD.

The eight distinct categories map to four monetary values; probabilities are supplied, not assumed equal.

\begin{center}\begin{tabular}{ccc}\toprule
Category & Given probability & $\GTenHLValue{X}$ (USD)\\\midrule
\GTenHLOutcome{No adjustment} & $\GTenHLProbability{0.07}=\GTenHLProbability{7/100}$ & $\GTenHLValue{0}$\\[4pt]
\GTenHLOutcome{Apology} & $\GTenHLProbability{0.09}=\GTenHLProbability{9/100}$ & $\GTenHLValue{0}$\\[4pt]
\GTenHLOutcome{Seat change} & $\GTenHLProbability{0.13}=\GTenHLProbability{13/100}$ & $\GTenHLValue{7}$\\[4pt]
\GTenHLOutcome{Programme} & $\GTenHLProbability{0.09}=\GTenHLProbability{9/100}$ & $\GTenHLValue{7}$\\[4pt]
\GTenHLOutcome{Partial refund} & $\GTenHLProbability{0.18}=\GTenHLProbability{9/50}$ & $\GTenHLValue{17}$\\[4pt]
\GTenHLOutcome{Replacement} & $\GTenHLProbability{0.22}=\GTenHLProbability{11/50}$ & $\GTenHLValue{17}$\\[4pt]
\GTenHLOutcome{Voucher} & $\GTenHLProbability{0.08}=\GTenHLProbability{2/25}$ & $\GTenHLValue{26}$\\[4pt]
\GTenHLOutcome{Cash} & $\GTenHLProbability{0.14}=\GTenHLProbability{7/50}$ & $\GTenHLValue{26}$\\[4pt]
\bottomrule\end{tabular}\end{center}

\par\vspace{3mm}\secInicio[][sec:q4]
\clearpage\textbf{Question 4 --- Task 1 continued: grouped \GTenHLProbability{PMF} and graph}


\begin{center}\begin{tabular}{@{}p{.30\linewidth}c p{.54\linewidth}@{}}\toprule
Complete outcomes & $\GTenHLValue{x}$ (USD) & Grouped \GTenHLProbability{PMF}\\\midrule
\GTenHLOutcome{No adjustment}, \GTenHLOutcome{Apology} & $\GTenHLValue{0}$ & $\begin{aligned}[t]p_X(\GTenHLValue{0})&=\GTenHLProbability{7/100}+\GTenHLProbability{9/100}\\&=\GTenHLProbability{4/25}=\GTenHLProbability{0.16}=\GTenHLProbability{16\%}\end{aligned}$\\[4pt]
\GTenHLOutcome{Seat change}, \GTenHLOutcome{Programme} & $\GTenHLValue{7}$ & $\begin{aligned}[t]p_X(\GTenHLValue{7})&=\GTenHLProbability{13/100}+\GTenHLProbability{9/100}\\&=\GTenHLProbability{11/50}=\GTenHLProbability{0.22}=\GTenHLProbability{22\%}\end{aligned}$\\[4pt]
\GTenHLOutcome{Partial refund}, \GTenHLOutcome{Replacement} & $\GTenHLValue{17}$ & $\begin{aligned}[t]p_X(\GTenHLValue{17})&=\GTenHLProbability{9/50}+\GTenHLProbability{11/50}\\&=\GTenHLProbability{2/5}=\GTenHLProbability{0.4}=\GTenHLProbability{40\%}\end{aligned}$\\[4pt]
\GTenHLOutcome{Voucher}, \GTenHLOutcome{Cash} & $\GTenHLValue{26}$ & $\begin{aligned}[t]p_X(\GTenHLValue{26})&=\GTenHLProbability{2/25}+\GTenHLProbability{7/50}\\&=\GTenHLProbability{11/50}=\GTenHLProbability{0.22}=\GTenHLProbability{22\%}\end{aligned}$\\[4pt]
\bottomrule\end{tabular}\end{center}
\begin{center}\begin{tikzpicture}\begin{axis}[keyaxis,width=.86\linewidth,height=5cm,xmin=-2,xmax=28,xtick={0,7,17,26},xticklabel style={rotate=45,anchor=north east},ymin=0,ymax=.7,ylabel={$\GTenHLProbability{p_X(x)}$}]
\draw[fill=keyMassFill,draw=keyMassColor] (axis cs:-0.225,0) rectangle (axis cs:0.225,4/25);
\draw[fill=keyMassFill,draw=keyMassColor] (axis cs:6.775,0) rectangle (axis cs:7.225,11/50);
\draw[fill=keyMassFill,draw=keyMassColor] (axis cs:16.775,0) rectangle (axis cs:17.225,2/5);
\draw[fill=keyMassFill,draw=keyMassColor] (axis cs:25.775,0) rectangle (axis cs:26.225,11/50);
\end{axis}\end{tikzpicture}\end{center}
\Needspace{5\baselineskip}\textbf{PMF checks.} All masses are positive. Normalization: $(4/25+11/50+2/5+11/50)=1$. Set $p_X(x)=0$ outside the support; gaps carry no mass.

\textbf{Scoring note.} These explanatory positivity and normalization checks carry no separate points; omitting them does not reduce the task score.


\par\smallskip
\noindent\begin{minipage}{\linewidth}\centering
\textbf{Proposed partial credit.}\par\smallskip
\begin{tabular}{@{}l r@{}}\toprule
\CreditDescription{\textbf{Assessed step or criterion}} & \textbf{Points}\\\midrule
\CreditDescription{Variable, units and support} & 0.125\\
\CreditDescription{Category mapping, grouping and the PMF table} & 0.25\\
\CreditDescription{The labeled, correctly spaced graph} & 0.125\\
\midrule
\CreditDescription{\textbf{Total}} & \textbf{0.5}\\
\bottomrule\end{tabular}
\end{minipage}\par\smallskip

\par\vspace{3mm}\secInicio[][sec:q4]

\clearpage
\textbf{Question 4 --- continued}\quad\hyperref[sec:q4]{Original description}

\Needspace{8\baselineskip}
\subsection*{Task 2 --- question and worked solution (proposed allocation) \score{1}}
\label{sol:q4t2}

\noindent\textbf{\boldmath Construct the complete CDF table, including both tails, and draw a labeled step graph. Use open and filled endpoints correctly and preserve actual monetary spacing.}\par\nopagebreak

For $\HLCumulativeProbability{F_X(x)}=P(X\leq x)$, add the masses in increasing monetary order.


\begin{center}\begin{tabular}{@{}l l c@{}}\toprule
Cutoff interval (USD) & Cumulative calculation & $\HLCumulativeProbability{F_X(x)}$\\\midrule
$x<\GTenHLValue{0}$ & No included support values & $\HLCumulativeProbability{0}$\\[4pt]
$\GTenHLValue{0}\leq x<\GTenHLValue{7}$ & $\begin{aligned}[t]&\GTenHLProbability{4/25}\end{aligned}$ & $\HLCumulativeProbability{4/25}$\\[4pt]
$\GTenHLValue{7}\leq x<\GTenHLValue{17}$ & $\begin{aligned}[t]&\GTenHLProbability{4/25}+\GTenHLProbability{11/50}\end{aligned}$ & $\HLCumulativeProbability{19/50}$\\[4pt]
$\GTenHLValue{17}\leq x<\GTenHLValue{26}$ & $\begin{aligned}[t]&\GTenHLProbability{4/25}+\GTenHLProbability{11/50}+\GTenHLProbability{2/5}\end{aligned}$ & $\HLCumulativeProbability{39/50}$\\[4pt]
$\GTenHLValue{26}\leq x$ & $\begin{aligned}[t]&\GTenHLProbability{4/25}+\GTenHLProbability{11/50}+\GTenHLProbability{2/5}+\GTenHLProbability{11/50}\end{aligned}$ & $\HLCumulativeProbability{1}$\\[4pt]
\bottomrule\end{tabular}\end{center}
\begin{center}\begin{tikzpicture}\begin{axis}[keyaxis,width=.86\linewidth,height=7cm,xmin=-2,xmax=28,xtick={0,7,17,26},xticklabel style={rotate=45,anchor=north east},ymin=-.07,ymax=1.1,ytick={0,.25,.5,.75,1},ylabel={$\HLCumulativeProbability{F_X(x)}$}]
\addplot[keyCDFColor,thick,<-] coordinates {(-2,0)(0,0)};
\addplot[keyCDFColor,thick] coordinates {(0,4/25)(7,4/25)};
\addplot[keyCDFColor,thick] coordinates {(7,19/50)(17,19/50)};
\addplot[keyCDFColor,thick] coordinates {(17,39/50)(26,39/50)};
\addplot[keyCDFColor,thick,->] coordinates {(26,1)(28,1)};
\addplot[keyCDFColor,only marks,mark=*,mark size=1pt,mark options={fill=white}] coordinates {(0,0)(7,4/25)(17,19/50)(26,39/50)};
\addplot[keyCDFColor,only marks,mark=*,mark size=1pt] coordinates {(0,4/25)(7,19/50)(17,39/50)(26,1)};
\end{axis}\end{tikzpicture}\end{center}
\par\vspace{3mm}\secInicio[][sec:q4]
\clearpage\textbf{Question 4 --- Task 2 continued: comparison}

\textbf{Same distribution: CDF above, PMF below.}
\begin{center}\begin{minipage}{\linewidth}\centering\begin{tikzpicture}
\pgfplotsset{keySupport/.style={xmin=-2,xmax=28,xtick={0,7,17,26},xticklabel style={rotate=45,anchor=north east}}}
\begin{axis}[keyCombined,keySupport,name=keyTop,anchor=south west,ylabel={$\HLCumulativeProbability{F_X(x)}$}]
\keyPhantom{0}{4/25}{0}
\keyPhantom{7}{11/50}{4/25}
\keyPhantom{17}{2/5}{19/50}
\keyPhantom{26}{11/50}{39/50}
\addplot[keyStep,<-] coordinates {(-2,0)(0,0)};
\addplot[keyStep] coordinates {(0,4/25)(7,4/25)};
\addplot[keyStep] coordinates {(7,19/50)(17,19/50)};
\addplot[keyStep] coordinates {(17,39/50)(26,39/50)};
\addplot[keyStep,->] coordinates {(26,1)(28,1)};
\addplot[keyExcluded] coordinates {(0,0)(7,4/25)(17,19/50)(26,39/50)};
\addplot[keyIncluded] coordinates {(0,4/25)(7,19/50)(17,39/50)(26,1)};
\end{axis}
\begin{axis}[keyCombined,keySupport,at={(keyTop.south west)},anchor=north west,yshift=-1.35cm,ylabel={$\GTenHLProbability{p_X(x)}$}]
\keyMass{0}{4/25}
\keyMass{7}{11/50}
\keyMass{17}{2/5}
\keyMass{26}{11/50}
\addlegendimage{keyStep}\addlegendentry{Top: CDF steps}
\addlegendimage{keyPhantom,mark=*}\addlegendentry{Top: PMF shadows (jumps)}
\addlegendimage{keyMass,mark=*}\addlegendentry{Bottom: PMF stems}
\addlegendimage{keyCorrespondence}\addlegendentry{Midpoint correspondence}
\end{axis}
\keyCorrespondence{0}
\keyCorrespondence{7}
\keyCorrespondence{17}
\keyCorrespondence{26}
\end{tikzpicture}\end{minipage}\end{center}\keyConnectionReading

\textbf{CDF checks.}\begin{itemize}\item The intervals cover every real cutoff; all heights lie in $[0,1]$ and never decrease.\item Each segment includes its left endpoint and excludes its right endpoint. Filled endpoints give the values at jumps, so the CDF is right-continuous.\item The lower-tail limit is 0 and the upper-tail limit is 1; arrows extend both tails.\item Successive differences recover the PMF: $4/25-0=4/25,\quad 19/50-4/25=11/50,\quad 39/50-19/50=2/5,\quad 1-39/50=11/50$. Thus $F_X(a)-F_X(a^-)=p_X(a)$; unsupported cutoffs have zero jump and zero mass.\end{itemize}


\par\smallskip
\noindent\begin{minipage}{\linewidth}\centering
\textbf{Proposed partial credit.}\par\smallskip
\begin{tabular}{@{}l r@{}}\toprule
\CreditDescription{\textbf{Assessed step or criterion}} & \textbf{Points}\\\midrule
\CreditDescription{Cumulative sums} & 0.25\\
\CreditDescription{The complete interval table and tails} & 0.25\\
\CreditDescription{The labeled, correctly spaced step graph and endpoints} & 0.25\\
\CreditDescription{The CDF and jump checks} & 0.25\\
\midrule
\CreditDescription{\textbf{Total}} & \textbf{1}\\
\bottomrule\end{tabular}
\end{minipage}\par\smallskip

\par\vspace{3mm}\secInicio[][sec:q4]

\clearpage
\textbf{Question 4 --- continued}\quad\hyperref[sec:q4]{Original description}

\Needspace{8\baselineskip}
\subsection*{Task 3 (c) --- question and worked solution (proposed allocation) \score{0.5}}
\label{sol:q4t3}

\noindent\textbf{\boldmath Calculate and interpret (i) the probability that the payout is less than 7 USD and (ii) the probability that it is at least 17 USD.}\par\nopagebreak

\textbf{Allocation:} (i) 0.25 points; (ii) 0.25 points.

\textbf{(i)} Include only $\GTenHLValue{0}$ USD; exclude the boundary mass at 7 USD.
\[\begin{aligned}P(\GTenHLValue{X}<\GTenHLValue{7})&=p_X(\GTenHLValue{0})\\&=\GTenHLProbability{4/25}\\&=\GTenHLProbability{4/25}=\GTenHLProbability{0.16}=\GTenHLProbability{16\%}.\end{aligned}\]
The payout for one case is below 7 USD with probability $\GTenHLProbability{4/25}$, approximately 16.00\%.

\textbf{(ii)} Include only $\GTenHLValue{17},\GTenHLValue{26}$ USD; include the boundary mass at 17 USD.
\[\begin{aligned}P(\GTenHLValue{X}\geq\GTenHLValue{17})&=p_X(\GTenHLValue{17})+p_X(\GTenHLValue{26})\\&=\GTenHLProbability{2/5}+\GTenHLProbability{11/50}\\&=\GTenHLProbability{31/50}=\GTenHLProbability{0.62}=\GTenHLProbability{62\%}.\end{aligned}\]
The payout for one case is at least 17 USD with probability $\GTenHLProbability{31/50}$, approximately 62.00\%.



\par\smallskip
\noindent\begin{minipage}{\linewidth}\centering
\textbf{Proposed partial credit.}\par\smallskip
\begin{tabular}{@{}l r@{}}\toprule
\CreditDescription{\textbf{Assessed step or criterion}} & \textbf{Points}\\\midrule
\CreditDescription{(i): selecting exactly the included PMF masses with the correct boundary} & 0.125\\
\CreditDescription{(i): summing them correctly} & 0.0625\\
\CreditDescription{(i): the contextual interpretation} & 0.0625\\
\CreditDescription{(ii): selecting exactly the included PMF masses with the correct boundary} & 0.125\\
\CreditDescription{(ii): summing them correctly} & 0.0625\\
\CreditDescription{(ii): the contextual interpretation} & 0.0625\\
\midrule
\CreditDescription{\textbf{Total}} & \textbf{0.5}\\
\bottomrule\end{tabular}
\end{minipage}\par\smallskip

Each part totals 0.25 points.

\par\vspace{3mm}\secInicio[][sec:q4]

\textbf{Question total: $0.5+1+0.5=2$ points for C2.}

\clearpage
\sectionScored{Courier service goodwill credits}{C2-\sml{4}}{sec:q5}

For one shipment, a courier company records exactly one of seven distinct, mutually exclusive and exhaustive service-resolution outcomes: A (parcel delivered on time); B (collection rescheduled); C (short transit delay); D (duplicate label corrected); E (alternative depot accepted); F (later delivery accepted); G (parcel returned by carrier). These named outcomes form the sample space. Outcomes A and B award 0 USD service credit; C and D award 7 USD; E, F and G award 15 USD. Thus seven outcomes map to three possible credit values. B is 4 times as likely as A; C is 8 times as likely as A; D is 4 times as likely as B; E is as likely as B; F is 4 times as likely as E; G is as likely as F. There are no other payments or adjustments. Determine the outcome probabilities from these likelihood relationships and record the service credit.

\Needspace{8\baselineskip}
\subsection*{Task 1 --- question and worked solution (proposed allocation) \score{1}}
\label{sol:q5t1}

\noindent\textbf{\boldmath Define the random variable and support. Derive the seven named outcome probabilities using normalization, map each outcome to its credit value, and construct the grouped PMF table and labeled graph.}\par\nopagebreak

Let $\GTenHLValue{X}$ be the service credit for one shipment, in \GTenHLValue{USD}. $\GTenHLValue{S_X}=\{\GTenHLValue{0},\GTenHLValue{7},\GTenHLValue{15}\}$ USD.

The sample space is $\Omega=\{A,B,C,D,E,F,G\}$. Set $\GTenHLParameter{a}=P(A)$, the baseline outcome probability, not the grouped mass $p_X(0)$.
\begin{align*}P(A)&=a,&P(B)&=4a,&P(C)&=8a,\\P(D)&=4P(B)=16a,&P(E)&=P(B)=4a,\\P(F)&=4P(E)=16a,&P(G)&=P(F)=16a.\end{align*}
Normalize over all seven exhaustive outcomes:
\[a+4a+8a+16a+4a+16a+16a=1,\qquad 65a=1,\qquad \GTenHLParameter{a}=\frac1{65}.\]

\begin{center}\begin{tabular}{ccc}\toprule
Named outcome & Outcome probability & $\GTenHLValue{X}$ (USD)\\\midrule
\GTenHLOutcome{A} & $1\GTenHLParameter{a}=\GTenHLProbability{1/65}$ & $\GTenHLValue{0}$\\[4pt]
\GTenHLOutcome{B} & $4\GTenHLParameter{a}=\GTenHLProbability{4/65}$ & $\GTenHLValue{0}$\\[4pt]
\GTenHLOutcome{C} & $8\GTenHLParameter{a}=\GTenHLProbability{8/65}$ & $\GTenHLValue{7}$\\[4pt]
\GTenHLOutcome{D} & $16\GTenHLParameter{a}=\GTenHLProbability{16/65}$ & $\GTenHLValue{7}$\\[4pt]
\GTenHLOutcome{E} & $4\GTenHLParameter{a}=\GTenHLProbability{4/65}$ & $\GTenHLValue{15}$\\[4pt]
\GTenHLOutcome{F} & $16\GTenHLParameter{a}=\GTenHLProbability{16/65}$ & $\GTenHLValue{15}$\\[4pt]
\GTenHLOutcome{G} & $16\GTenHLParameter{a}=\GTenHLProbability{16/65}$ & $\GTenHLValue{15}$\\[4pt]
\bottomrule\end{tabular}\end{center}

\par\vspace{3mm}\secInicio[][sec:q5]
\clearpage\textbf{Question 5 --- Task 1 continued: grouped \GTenHLProbability{PMF} and graph}


\begin{center}\begin{tabular}{@{}p{.30\linewidth}c p{.54\linewidth}@{}}\toprule
Complete outcomes & $\GTenHLValue{x}$ (USD) & Grouped \GTenHLProbability{PMF}\\\midrule
\GTenHLOutcome{A}, \GTenHLOutcome{B} & $\GTenHLValue{0}$ & $\begin{aligned}[t]p_X(\GTenHLValue{0})&=\GTenHLProbability{1/65}+\GTenHLProbability{4/65}\\&=\GTenHLProbability{1/13}\approx\GTenHLProbability{0.0769}=\GTenHLProbability{7.69\%}\end{aligned}$\\[4pt]
\GTenHLOutcome{C}, \GTenHLOutcome{D} & $\GTenHLValue{7}$ & $\begin{aligned}[t]p_X(\GTenHLValue{7})&=\GTenHLProbability{8/65}+\GTenHLProbability{16/65}\\&=\GTenHLProbability{24/65}\approx\GTenHLProbability{0.3692}=\GTenHLProbability{36.92\%}\end{aligned}$\\[4pt]
\GTenHLOutcome{E}, \GTenHLOutcome{F}, \GTenHLOutcome{G} & $\GTenHLValue{15}$ & $\begin{aligned}[t]p_X(\GTenHLValue{15})&=\GTenHLProbability{4/65}+\GTenHLProbability{16/65}+\GTenHLProbability{16/65}\\&=\GTenHLProbability{36/65}\approx\GTenHLProbability{0.5538}=\GTenHLProbability{55.38\%}\end{aligned}$\\[4pt]
\bottomrule\end{tabular}\end{center}
\begin{center}\begin{tikzpicture}\begin{axis}[keyaxis,width=.86\linewidth,height=5cm,xmin=-2,xmax=17,xtick={0,7,15},xticklabel style={rotate=45,anchor=north east},ymin=0,ymax=.7,ylabel={$\GTenHLProbability{p_X(x)}$}]
\draw[fill=keyMassFill,draw=keyMassColor] (axis cs:-0.225,0) rectangle (axis cs:0.225,1/13);
\draw[fill=keyMassFill,draw=keyMassColor] (axis cs:6.775,0) rectangle (axis cs:7.225,24/65);
\draw[fill=keyMassFill,draw=keyMassColor] (axis cs:14.775,0) rectangle (axis cs:15.225,36/65);
\end{axis}\end{tikzpicture}\end{center}
\Needspace{5\baselineskip}\textbf{PMF checks.} All masses are positive. Normalization: $(1/13+24/65+36/65)=1$. Set $p_X(x)=0$ outside the support; gaps carry no mass.

\textbf{Scoring note.} These explanatory positivity and normalization checks carry no separate points; omitting them does not reduce the task score.
\textbf{Relationship check.} The outcome numerators satisfy $4=4(1),\ 8=8(1),\ 16=4(4),\ 4=4,\ 16=4(4),\ 16=16$, reproducing all six relationships.


\par\smallskip
\noindent\begin{minipage}{\linewidth}\centering
\textbf{Proposed partial credit.}\par\smallskip
\begin{tabular}{@{}l r@{}}\toprule
\CreditDescription{\textbf{Assessed step or criterion}} & \textbf{Points}\\\midrule
\CreditDescription{Variable, units and support} & 0.125\\
\CreditDescription{Translating the relationships and solving normalization} & 0.375\\
\CreditDescription{All seven outcome probabilities, mappings and three grouped PMF entries} & 0.25\\
\CreditDescription{The labeled, correctly spaced graph} & 0.125\\
\CreditDescription{Relationship checks} & 0.125\\
\midrule
\CreditDescription{\textbf{Total}} & \textbf{1}\\
\bottomrule\end{tabular}
\end{minipage}\par\smallskip

\par\vspace{3mm}\secInicio[][sec:q5]

\clearpage
\textbf{Question 5 --- continued}\quad\hyperref[sec:q5]{Original description}

\Needspace{8\baselineskip}
\subsection*{Task 2 --- question and worked solution (proposed allocation) \score{2}}
\label{sol:q5t2}

\noindent\textbf{\boldmath Construct the complete CDF table, including both tails, and draw a labeled step graph. Use open and filled endpoints correctly and preserve actual monetary spacing.}\par\nopagebreak

For $\HLCumulativeProbability{F_X(x)}=P(X\leq x)$, add the masses in increasing monetary order.


\begin{center}\begin{tabular}{@{}l l c@{}}\toprule
Cutoff interval (USD) & Cumulative calculation & $\HLCumulativeProbability{F_X(x)}$\\\midrule
$x<\GTenHLValue{0}$ & No included support values & $\HLCumulativeProbability{0}$\\[4pt]
$\GTenHLValue{0}\leq x<\GTenHLValue{7}$ & $\begin{aligned}[t]&\GTenHLProbability{1/13}\end{aligned}$ & $\HLCumulativeProbability{1/13}$\\[4pt]
$\GTenHLValue{7}\leq x<\GTenHLValue{15}$ & $\begin{aligned}[t]&\GTenHLProbability{1/13}+\GTenHLProbability{24/65}\end{aligned}$ & $\HLCumulativeProbability{29/65}$\\[4pt]
$\GTenHLValue{15}\leq x$ & $\begin{aligned}[t]&\GTenHLProbability{1/13}+\GTenHLProbability{24/65}+\GTenHLProbability{36/65}\end{aligned}$ & $\HLCumulativeProbability{1}$\\[4pt]
\bottomrule\end{tabular}\end{center}
\begin{center}\begin{tikzpicture}\begin{axis}[keyaxis,width=.86\linewidth,height=7cm,xmin=-2,xmax=17,xtick={0,7,15},xticklabel style={rotate=45,anchor=north east},ymin=-.07,ymax=1.1,ytick={0,.25,.5,.75,1},ylabel={$\HLCumulativeProbability{F_X(x)}$}]
\addplot[keyCDFColor,thick,<-] coordinates {(-2,0)(0,0)};
\addplot[keyCDFColor,thick] coordinates {(0,1/13)(7,1/13)};
\addplot[keyCDFColor,thick] coordinates {(7,29/65)(15,29/65)};
\addplot[keyCDFColor,thick,->] coordinates {(15,1)(17,1)};
\addplot[keyCDFColor,only marks,mark=*,mark size=1pt,mark options={fill=white}] coordinates {(0,0)(7,1/13)(15,29/65)};
\addplot[keyCDFColor,only marks,mark=*,mark size=1pt] coordinates {(0,1/13)(7,29/65)(15,1)};
\end{axis}\end{tikzpicture}\end{center}
\par\vspace{3mm}\secInicio[][sec:q5]
\clearpage\textbf{Question 5 --- Task 2 continued: comparison}

\textbf{Same distribution: CDF above, PMF below.}
\begin{center}\begin{minipage}{\linewidth}\centering\begin{tikzpicture}
\pgfplotsset{keySupport/.style={xmin=-2,xmax=17,xtick={0,7,15},xticklabel style={rotate=45,anchor=north east}}}
\begin{axis}[keyCombined,keySupport,name=keyTop,anchor=south west,ylabel={$\HLCumulativeProbability{F_X(x)}$}]
\keyPhantom{0}{1/13}{0}
\keyPhantom{7}{24/65}{1/13}
\keyPhantom{15}{36/65}{29/65}
\addplot[keyStep,<-] coordinates {(-2,0)(0,0)};
\addplot[keyStep] coordinates {(0,1/13)(7,1/13)};
\addplot[keyStep] coordinates {(7,29/65)(15,29/65)};
\addplot[keyStep,->] coordinates {(15,1)(17,1)};
\addplot[keyExcluded] coordinates {(0,0)(7,1/13)(15,29/65)};
\addplot[keyIncluded] coordinates {(0,1/13)(7,29/65)(15,1)};
\end{axis}
\begin{axis}[keyCombined,keySupport,at={(keyTop.south west)},anchor=north west,yshift=-1.35cm,ylabel={$\GTenHLProbability{p_X(x)}$}]
\keyMass{0}{1/13}
\keyMass{7}{24/65}
\keyMass{15}{36/65}
\addlegendimage{keyStep}\addlegendentry{Top: CDF steps}
\addlegendimage{keyPhantom,mark=*}\addlegendentry{Top: PMF shadows (jumps)}
\addlegendimage{keyMass,mark=*}\addlegendentry{Bottom: PMF stems}
\addlegendimage{keyCorrespondence}\addlegendentry{Midpoint correspondence}
\end{axis}
\keyCorrespondence{0}
\keyCorrespondence{7}
\keyCorrespondence{15}
\end{tikzpicture}\end{minipage}\end{center}\keyConnectionReading

\textbf{CDF checks.}\begin{itemize}\item The intervals cover every real cutoff; all heights lie in $[0,1]$ and never decrease.\item Each segment includes its left endpoint and excludes its right endpoint. Filled endpoints give the values at jumps, so the CDF is right-continuous.\item The lower-tail limit is 0 and the upper-tail limit is 1; arrows extend both tails.\item Successive differences recover the PMF: $1/13-0=1/13,\quad 29/65-1/13=24/65,\quad 1-29/65=36/65$. Thus $F_X(a)-F_X(a^-)=p_X(a)$; unsupported cutoffs have zero jump and zero mass.\end{itemize}


\par\smallskip
\noindent\begin{minipage}{\linewidth}\centering
\textbf{Proposed partial credit.}\par\smallskip
\begin{tabular}{@{}l r@{}}\toprule
\CreditDescription{\textbf{Assessed step or criterion}} & \textbf{Points}\\\midrule
\CreditDescription{Cumulative sums} & 0.5\\
\CreditDescription{The complete interval table and tails} & 0.5\\
\CreditDescription{The labeled, correctly spaced step graph and endpoints} & 0.5\\
\CreditDescription{The CDF and jump checks} & 0.5\\
\midrule
\CreditDescription{\textbf{Total}} & \textbf{2}\\
\bottomrule\end{tabular}
\end{minipage}\par\smallskip

\par\vspace{3mm}\secInicio[][sec:q5]

\clearpage
\textbf{Question 5 --- continued}\quad\hyperref[sec:q5]{Original description}

\Needspace{8\baselineskip}
\subsection*{Task 3 (c) --- question and worked solution (proposed allocation) \score{1}}
\label{sol:q5t3}

\noindent\textbf{\boldmath Calculate and interpret (i) the probability that the service credit is less than 7 USD and (ii) the probability that it is at least 15 USD.}\par\nopagebreak

\textbf{Allocation:} (i) 0.5 points; (ii) 0.5 points.

\textbf{(i)} Include only $\GTenHLValue{0}$ USD; exclude the boundary mass at 7 USD.
\[\begin{aligned}P(\GTenHLValue{X}<\GTenHLValue{7})&=p_X(\GTenHLValue{0})\\&=\GTenHLProbability{1/13}\\&=\GTenHLProbability{1/13}\approx\GTenHLProbability{0.0769}=\GTenHLProbability{7.69\%}.\end{aligned}\]
The service credit for one case is below 7 USD with probability $\GTenHLProbability{1/13}$, approximately 7.69\%.

\textbf{(ii)} Include only $\GTenHLValue{15}$ USD; include the boundary mass at 15 USD.
\[\begin{aligned}P(\GTenHLValue{X}\geq\GTenHLValue{15})&=p_X(\GTenHLValue{15})\\&=\GTenHLProbability{36/65}\\&=\GTenHLProbability{36/65}\approx\GTenHLProbability{0.5538}=\GTenHLProbability{55.38\%}.\end{aligned}\]
The service credit for one case is at least 15 USD with probability $\GTenHLProbability{36/65}$, approximately 55.38\%.



\par\smallskip
\noindent\begin{minipage}{\linewidth}\centering
\textbf{Proposed partial credit.}\par\smallskip
\begin{tabular}{@{}l r@{}}\toprule
\CreditDescription{\textbf{Assessed step or criterion}} & \textbf{Points}\\\midrule
\CreditDescription{(i): selecting exactly the included PMF masses with the correct boundary} & 0.25\\
\CreditDescription{(i): summing them correctly} & 0.125\\
\CreditDescription{(i): the contextual interpretation} & 0.125\\
\CreditDescription{(ii): selecting exactly the included PMF masses with the correct boundary} & 0.25\\
\CreditDescription{(ii): summing them correctly} & 0.125\\
\CreditDescription{(ii): the contextual interpretation} & 0.125\\
\midrule
\CreditDescription{\textbf{Total}} & \textbf{1}\\
\bottomrule\end{tabular}
\end{minipage}\par\smallskip

Each part totals 0.5 points.

\par\vspace{3mm}\secInicio[][sec:q5]

\textbf{Question total: $1+2+1=4$ points for C2.}

\clearpage
\sectionScored{Animal shelter raffle coupons}{C2-\sml{6}}{sec:q6}

An animal shelter raffle has 10 physically distinct prize coupons: 5 worth 0 USD each, 3 worth 5 USD each and 2 worth 12 USD each. Choose a subset of three uniformly without replacement, so every subset of three physically distinct items has the same chance. Ignore selection order. First add the face values. Then add one bonus of 7 USD if the selected subset contains at least one 5 USD item and at least one 12 USD item; otherwise add no bonus. Finally, if this adjusted amount exceeds 27 USD, pay only 27 USD; otherwise pay the adjusted amount. No other payments or deductions apply. Record the final payout.

\Needspace{8\baselineskip}
\subsection*{Task 1 --- question and worked solution (proposed allocation) \score{1.5}}
\label{sol:q6t1}

\noindent\textbf{\boldmath Define the random variable and support. Count the distinct-ticket selections, apply the bonus and then the cap, aggregate equal final payouts, and derive the PMF. Produce its table and labeled graph.}\par\nopagebreak

Let $\GTenHLValue{X}$ be the final payout for one three-item subset, in \GTenHLValue{USD}. $\GTenHLValue{S_X}=\{\GTenHLValue{0},\GTenHLValue{5},\GTenHLValue{10},\GTenHLValue{12},\GTenHLValue{15},\GTenHLValue{24},\GTenHLValue{27}\}$ USD.

Label the distinct items $Z_1,\ldots,Z_5$, $L_1,\ldots,L_3$ and $H_1,H_2$, worth 0, 5 and 12 USD respectively. There are $\binom{10}3=120$ equally likely unordered subsets, each with probability $1/120$.

A composition $(a,b,c)$ counts zero-, 5- and 12-dollar items. It satisfies $a+b+c=3$, $0\leq a\leq 5$, $0\leq b\leq 3$ and $0\leq c\leq 2$. Its count is $N(a,b,c)=\binom{5}a\binom{3}b\binom{2}c$ and its probability is $N(a,b,c)/120$. Let $I=1$ when $b\geq1$ and $c\geq1$, and 0 otherwise. First calculate $Y=5b+12c+7I$, then $X=\min\{Y,27\}$: add the bonus once before the cap.

\begin{center}\begin{tabular}{ccccccc}\toprule
$(a,b,c)$ & Subset count & Probability & Face sum & Bonus & $Y$ & $X$\\\midrule
$\GTenHLOutcome{(3,0,0)}$ & $\binom{5}{3}\binom{3}{0}\binom{2}{0}=10$ & $\GTenHLProbability{1/12}$ & $\GTenHLValue{0}$ & $\GTenHLValue{0}$ & $\GTenHLSecond{0}$ & $\GTenHLValue{0}$\\[4pt]
$\GTenHLOutcome{(2,1,0)}$ & $\binom{5}{2}\binom{3}{1}\binom{2}{0}=30$ & $\GTenHLProbability{1/4}$ & $\GTenHLValue{5}$ & $\GTenHLValue{0}$ & $\GTenHLSecond{5}$ & $\GTenHLValue{5}$\\[4pt]
$\GTenHLOutcome{(1,2,0)}$ & $\binom{5}{1}\binom{3}{2}\binom{2}{0}=15$ & $\GTenHLProbability{1/8}$ & $\GTenHLValue{10}$ & $\GTenHLValue{0}$ & $\GTenHLSecond{10}$ & $\GTenHLValue{10}$\\[4pt]
$\GTenHLOutcome{(0,3,0)}$ & $\binom{5}{0}\binom{3}{3}\binom{2}{0}=1$ & $\GTenHLProbability{1/120}$ & $\GTenHLValue{15}$ & $\GTenHLValue{0}$ & $\GTenHLSecond{15}$ & $\GTenHLValue{15}$\\[4pt]
$\GTenHLOutcome{(2,0,1)}$ & $\binom{5}{2}\binom{3}{0}\binom{2}{1}=20$ & $\GTenHLProbability{1/6}$ & $\GTenHLValue{12}$ & $\GTenHLValue{0}$ & $\GTenHLSecond{12}$ & $\GTenHLValue{12}$\\[4pt]
$\GTenHLOutcome{(1,1,1)}$ & $\binom{5}{1}\binom{3}{1}\binom{2}{1}=30$ & $\GTenHLProbability{1/4}$ & $\GTenHLValue{17}$ & $\GTenHLValue{7}$ & $\GTenHLSecond{24}$ & $\GTenHLValue{24}$\\[4pt]
$\GTenHLOutcome{(0,2,1)}$ & $\binom{5}{0}\binom{3}{2}\binom{2}{1}=6$ & $\GTenHLProbability{1/20}$ & $\GTenHLValue{22}$ & $\GTenHLValue{7}$ & $\GTenHLSecond{29}$ & $\GTenHLValue{27}$\\[4pt]
$\GTenHLOutcome{(1,0,2)}$ & $\binom{5}{1}\binom{3}{0}\binom{2}{2}=5$ & $\GTenHLProbability{1/24}$ & $\GTenHLValue{24}$ & $\GTenHLValue{0}$ & $\GTenHLSecond{24}$ & $\GTenHLValue{24}$\\[4pt]
$\GTenHLOutcome{(0,1,2)}$ & $\binom{5}{0}\binom{3}{1}\binom{2}{2}=3$ & $\GTenHLProbability{1/40}$ & $\GTenHLValue{29}$ & $\GTenHLValue{7}$ & $\GTenHLSecond{36}$ & $\GTenHLValue{27}$\\[4pt]
\bottomrule\end{tabular}\end{center}
All monetary columns are in USD. Listing $c=0,1,2$ and all feasible $b$ gives the nine compositions. Counts sum to $10+30+15+1+20+30+6+5+3=120$. Compositions are not equally likely: $(1,1,1)$ represents 30 subsets, while $(3,0,0)$ represents 10.

\textbf{Aggregation.} Below the cap, 24 USD combines $(1,1,1)$ after its bonus with $(1,0,2)$ without a bonus: $(30+5)/120=7/24$. At the cap, adjusted amounts 29 and 36 USD from $(0,2,1)$ and $(0,1,2)$ combine: $(6+3)/120=3/40$. Every other payout comes from one composition; 24 USD stays below the cap.

\par\vspace{3mm}\secInicio[][sec:q6]
\clearpage\textbf{Question 6 --- Task 1 continued: grouped \GTenHLProbability{PMF} and graph}


\begin{center}\begin{tabular}{@{}p{.30\linewidth}c p{.54\linewidth}@{}}\toprule
Complete outcomes & $\GTenHLValue{x}$ (USD) & Grouped \GTenHLProbability{PMF}\\\midrule
$\GTenHLOutcome{(3,0,0)}$ & $\GTenHLValue{0}$ & $\begin{aligned}[t]p_X(\GTenHLValue{0})&=\GTenHLProbability{1/12}\\&=\GTenHLProbability{1/12}\approx\GTenHLProbability{0.0833}=\GTenHLProbability{8.33\%}\end{aligned}$\\[4pt]
$\GTenHLOutcome{(2,1,0)}$ & $\GTenHLValue{5}$ & $\begin{aligned}[t]p_X(\GTenHLValue{5})&=\GTenHLProbability{1/4}\\&=\GTenHLProbability{1/4}=\GTenHLProbability{0.25}=\GTenHLProbability{25\%}\end{aligned}$\\[4pt]
$\GTenHLOutcome{(1,2,0)}$ & $\GTenHLValue{10}$ & $\begin{aligned}[t]p_X(\GTenHLValue{10})&=\GTenHLProbability{1/8}\\&=\GTenHLProbability{1/8}=\GTenHLProbability{0.125}=\GTenHLProbability{12.5\%}\end{aligned}$\\[4pt]
$\GTenHLOutcome{(2,0,1)}$ & $\GTenHLValue{12}$ & $\begin{aligned}[t]p_X(\GTenHLValue{12})&=\GTenHLProbability{1/6}\\&=\GTenHLProbability{1/6}\approx\GTenHLProbability{0.1667}=\GTenHLProbability{16.67\%}\end{aligned}$\\[4pt]
$\GTenHLOutcome{(0,3,0)}$ & $\GTenHLValue{15}$ & $\begin{aligned}[t]p_X(\GTenHLValue{15})&=\GTenHLProbability{1/120}\\&=\GTenHLProbability{1/120}\approx\GTenHLProbability{0.0083}=\GTenHLProbability{0.83\%}\end{aligned}$\\[4pt]
$\GTenHLOutcome{(1,1,1)}$, $\GTenHLOutcome{(1,0,2)}$ & $\GTenHLValue{24}$ & $\begin{aligned}[t]p_X(\GTenHLValue{24})&=\GTenHLProbability{1/4}+\GTenHLProbability{1/24}\\&=\GTenHLProbability{7/24}\approx\GTenHLProbability{0.2917}=\GTenHLProbability{29.17\%}\end{aligned}$\\[4pt]
$\GTenHLOutcome{(0,2,1)}$, $\GTenHLOutcome{(0,1,2)}$ & $\GTenHLValue{27}$ & $\begin{aligned}[t]p_X(\GTenHLValue{27})&=\GTenHLProbability{1/20}+\GTenHLProbability{1/40}\\&=\GTenHLProbability{3/40}=\GTenHLProbability{0.075}=\GTenHLProbability{7.5\%}\end{aligned}$\\[4pt]
\bottomrule\end{tabular}\end{center}
\begin{center}\begin{tikzpicture}\begin{axis}[keyaxis,width=.86\linewidth,height=5cm,xmin=-2,xmax=29,xtick={0,5,10,12,15,24,27},xticklabel style={rotate=45,anchor=north east},ymin=0,ymax=.7,ylabel={$\GTenHLProbability{p_X(x)}$}]
\draw[fill=keyMassFill,draw=keyMassColor] (axis cs:-0.225,0) rectangle (axis cs:0.225,1/12);
\draw[fill=keyMassFill,draw=keyMassColor] (axis cs:4.775,0) rectangle (axis cs:5.225,1/4);
\draw[fill=keyMassFill,draw=keyMassColor] (axis cs:9.775,0) rectangle (axis cs:10.225,1/8);
\draw[fill=keyMassFill,draw=keyMassColor] (axis cs:11.775,0) rectangle (axis cs:12.225,1/6);
\draw[fill=keyMassFill,draw=keyMassColor] (axis cs:14.775,0) rectangle (axis cs:15.225,1/120);
\draw[fill=keyMassFill,draw=keyMassColor] (axis cs:23.775,0) rectangle (axis cs:24.225,7/24);
\draw[fill=keyMassFill,draw=keyMassColor] (axis cs:26.775,0) rectangle (axis cs:27.225,3/40);
\end{axis}\end{tikzpicture}\end{center}
\Needspace{5\baselineskip}\textbf{PMF checks.} All masses are positive. Normalization: $(1/12+1/4+1/8+1/6+1/120+7/24+3/40)=1$. Set $p_X(x)=0$ outside the support; gaps carry no mass.

\textbf{Scoring note.} These explanatory positivity and normalization checks carry no separate points; omitting them does not reduce the task score.


\par\smallskip
\noindent\begin{minipage}{\linewidth}\centering
\textbf{Proposed partial credit.}\par\smallskip
\begin{tabular}{@{}l r@{}}\toprule
\CreditDescription{\textbf{Assessed step or criterion}} & \textbf{Points}\\\midrule
\CreditDescription{Variable, units and support} & 0.125\\
\CreditDescription{The denominator and exhaustive composition counts} & 0.5\\
\CreditDescription{The bonus, cap and aggregation} & 0.375\\
\CreditDescription{The grouped PMF table} & 0.25\\
\CreditDescription{The labeled, correctly spaced graph} & 0.25\\
\midrule
\CreditDescription{\textbf{Total}} & \textbf{1.5}\\
\bottomrule\end{tabular}
\end{minipage}\par\smallskip

\par\vspace{3mm}\secInicio[][sec:q6]

\clearpage
\textbf{Question 6 --- continued}\quad\hyperref[sec:q6]{Original description}

\Needspace{8\baselineskip}
\subsection*{Task 2 --- question and worked solution (proposed allocation) \score{3}}
\label{sol:q6t2}

\noindent\textbf{\boldmath Construct the complete CDF table, including both tails, and draw a labeled step graph. Use open and filled endpoints correctly and preserve actual monetary spacing.}\par\nopagebreak

For $\HLCumulativeProbability{F_X(x)}=P(X\leq x)$, add the masses in increasing monetary order.


\begin{center}\begin{tabular}{@{}l l c@{}}\toprule
Cutoff interval (USD) & Cumulative calculation & $\HLCumulativeProbability{F_X(x)}$\\\midrule
$x<\GTenHLValue{0}$ & No included support values & $\HLCumulativeProbability{0}$\\[4pt]
$\GTenHLValue{0}\leq x<\GTenHLValue{5}$ & $\begin{aligned}[t]&\GTenHLProbability{1/12}\end{aligned}$ & $\HLCumulativeProbability{1/12}$\\[4pt]
$\GTenHLValue{5}\leq x<\GTenHLValue{10}$ & $\begin{aligned}[t]&\GTenHLProbability{1/12}+\GTenHLProbability{1/4}\end{aligned}$ & $\HLCumulativeProbability{1/3}$\\[4pt]
$\GTenHLValue{10}\leq x<\GTenHLValue{12}$ & $\begin{aligned}[t]&\GTenHLProbability{1/12}+\GTenHLProbability{1/4}+\GTenHLProbability{1/8}\end{aligned}$ & $\HLCumulativeProbability{11/24}$\\[4pt]
$\GTenHLValue{12}\leq x<\GTenHLValue{15}$ & $\begin{aligned}[t]&\GTenHLProbability{1/12}+\GTenHLProbability{1/4}+\GTenHLProbability{1/8}+\GTenHLProbability{1/6}\end{aligned}$ & $\HLCumulativeProbability{5/8}$\\[4pt]
$\GTenHLValue{15}\leq x<\GTenHLValue{24}$ & $\begin{aligned}[t]&\GTenHLProbability{1/12}+\GTenHLProbability{1/4}+\GTenHLProbability{1/8}+\GTenHLProbability{1/6}\\&{}+\GTenHLProbability{1/120}\end{aligned}$ & $\HLCumulativeProbability{19/30}$\\[4pt]
$\GTenHLValue{24}\leq x<\GTenHLValue{27}$ & $\begin{aligned}[t]&\GTenHLProbability{1/12}+\GTenHLProbability{1/4}+\GTenHLProbability{1/8}+\GTenHLProbability{1/6}\\&{}+\GTenHLProbability{1/120}+\GTenHLProbability{7/24}\end{aligned}$ & $\HLCumulativeProbability{37/40}$\\[4pt]
$\GTenHLValue{27}\leq x$ & $\begin{aligned}[t]&\GTenHLProbability{1/12}+\GTenHLProbability{1/4}+\GTenHLProbability{1/8}+\GTenHLProbability{1/6}\\&{}+\GTenHLProbability{1/120}+\GTenHLProbability{7/24}+\GTenHLProbability{3/40}\end{aligned}$ & $\HLCumulativeProbability{1}$\\[4pt]
\bottomrule\end{tabular}\end{center}
\begin{center}\begin{tikzpicture}\begin{axis}[keyaxis,width=.86\linewidth,height=7cm,xmin=-2,xmax=29,xtick={0,5,10,12,15,24,27},xticklabel style={rotate=45,anchor=north east},ymin=-.07,ymax=1.1,ytick={0,.25,.5,.75,1},ylabel={$\HLCumulativeProbability{F_X(x)}$}]
\addplot[keyCDFColor,thick,<-] coordinates {(-2,0)(0,0)};
\addplot[keyCDFColor,thick] coordinates {(0,1/12)(5,1/12)};
\addplot[keyCDFColor,thick] coordinates {(5,1/3)(10,1/3)};
\addplot[keyCDFColor,thick] coordinates {(10,11/24)(12,11/24)};
\addplot[keyCDFColor,thick] coordinates {(12,5/8)(15,5/8)};
\addplot[keyCDFColor,thick] coordinates {(15,19/30)(24,19/30)};
\addplot[keyCDFColor,thick] coordinates {(24,37/40)(27,37/40)};
\addplot[keyCDFColor,thick,->] coordinates {(27,1)(29,1)};
\addplot[keyCDFColor,only marks,mark=*,mark size=1pt,mark options={fill=white}] coordinates {(0,0)(5,1/12)(10,1/3)(12,11/24)(15,5/8)(24,19/30)(27,37/40)};
\addplot[keyCDFColor,only marks,mark=*,mark size=1pt] coordinates {(0,1/12)(5,1/3)(10,11/24)(12,5/8)(15,19/30)(24,37/40)(27,1)};
\end{axis}\end{tikzpicture}\end{center}
\par\vspace{3mm}\secInicio[][sec:q6]
\clearpage\textbf{Question 6 --- Task 2 continued: comparison}

\textbf{Same distribution: CDF above, PMF below.}
\begin{center}\begin{minipage}{\linewidth}\centering\begin{tikzpicture}
\pgfplotsset{keySupport/.style={xmin=-2,xmax=29,xtick={0,5,10,12,15,24,27},xticklabel style={rotate=45,anchor=north east}}}
\begin{axis}[keyCombined,keySupport,name=keyTop,anchor=south west,ylabel={$\HLCumulativeProbability{F_X(x)}$}]
\keyPhantom{0}{1/12}{0}
\keyPhantom{5}{1/4}{1/12}
\keyPhantom{10}{1/8}{1/3}
\keyPhantom{12}{1/6}{11/24}
\keyPhantom{15}{1/120}{5/8}
\keyPhantom{24}{7/24}{19/30}
\keyPhantom{27}{3/40}{37/40}
\addplot[keyStep,<-] coordinates {(-2,0)(0,0)};
\addplot[keyStep] coordinates {(0,1/12)(5,1/12)};
\addplot[keyStep] coordinates {(5,1/3)(10,1/3)};
\addplot[keyStep] coordinates {(10,11/24)(12,11/24)};
\addplot[keyStep] coordinates {(12,5/8)(15,5/8)};
\addplot[keyStep] coordinates {(15,19/30)(24,19/30)};
\addplot[keyStep] coordinates {(24,37/40)(27,37/40)};
\addplot[keyStep,->] coordinates {(27,1)(29,1)};
\addplot[keyExcluded] coordinates {(0,0)(5,1/12)(10,1/3)(12,11/24)(15,5/8)(24,19/30)(27,37/40)};
\addplot[keyIncluded] coordinates {(0,1/12)(5,1/3)(10,11/24)(12,5/8)(15,19/30)(24,37/40)(27,1)};
\end{axis}
\begin{axis}[keyCombined,keySupport,at={(keyTop.south west)},anchor=north west,yshift=-1.35cm,ylabel={$\GTenHLProbability{p_X(x)}$}]
\keyMass{0}{1/12}
\keyMass{5}{1/4}
\keyMass{10}{1/8}
\keyMass{12}{1/6}
\keyMass{15}{1/120}
\keyMass{24}{7/24}
\keyMass{27}{3/40}
\addlegendimage{keyStep}\addlegendentry{Top: CDF steps}
\addlegendimage{keyPhantom,mark=*}\addlegendentry{Top: PMF shadows (jumps)}
\addlegendimage{keyMass,mark=*}\addlegendentry{Bottom: PMF stems}
\addlegendimage{keyCorrespondence}\addlegendentry{Midpoint correspondence}
\end{axis}
\keyCorrespondence{0}
\keyCorrespondence{5}
\keyCorrespondence{10}
\keyCorrespondence{12}
\keyCorrespondence{15}
\keyCorrespondence{24}
\keyCorrespondence{27}
\end{tikzpicture}\end{minipage}\end{center}\keyConnectionReading

\textbf{CDF checks.}\begin{itemize}\item The intervals cover every real cutoff; all heights lie in $[0,1]$ and never decrease.\item Each segment includes its left endpoint and excludes its right endpoint. Filled endpoints give the values at jumps, so the CDF is right-continuous.\item The lower-tail limit is 0 and the upper-tail limit is 1; arrows extend both tails.\item Successive differences recover the PMF: $1/12-0=1/12,\quad 1/3-1/12=1/4,\quad 11/24-1/3=1/8,\quad 5/8-11/24=1/6,\quad 19/30-5/8=1/120,\quad 37/40-19/30=7/24,\quad 1-37/40=3/40$. Thus $F_X(a)-F_X(a^-)=p_X(a)$; unsupported cutoffs have zero jump and zero mass.\end{itemize}


\par\smallskip
\noindent\begin{minipage}{\linewidth}\centering
\textbf{Proposed partial credit.}\par\smallskip
\begin{tabular}{@{}l r@{}}\toprule
\CreditDescription{\textbf{Assessed step or criterion}} & \textbf{Points}\\\midrule
\CreditDescription{Cumulative sums} & 0.75\\
\CreditDescription{The complete interval table and tails} & 0.75\\
\CreditDescription{The labeled, correctly spaced step graph and endpoints} & 0.75\\
\CreditDescription{The CDF and jump checks} & 0.75\\
\midrule
\CreditDescription{\textbf{Total}} & \textbf{3}\\
\bottomrule\end{tabular}
\end{minipage}\par\smallskip

\par\vspace{3mm}\secInicio[][sec:q6]

\clearpage
\textbf{Question 6 --- continued}\quad\hyperref[sec:q6]{Original description}

\Needspace{8\baselineskip}
\subsection*{Task 3 (c) --- question and worked solution (proposed allocation) \score{1.5}}
\label{sol:q6t3}

\noindent\textbf{\boldmath Calculate and interpret (i) the probability that the final payout is less than 12 USD and (ii) the probability that it is at least 24 USD.}\par\nopagebreak

\textbf{Allocation:} (i) 0.75 points; (ii) 0.75 points.

\textbf{(i)} Include only $\GTenHLValue{0},\GTenHLValue{5},\GTenHLValue{10}$ USD; exclude the boundary mass at 12 USD.
\[\begin{aligned}P(\GTenHLValue{X}<\GTenHLValue{12})&=p_X(\GTenHLValue{0})+p_X(\GTenHLValue{5})+p_X(\GTenHLValue{10})\\&=\GTenHLProbability{1/12}+\GTenHLProbability{1/4}+\GTenHLProbability{1/8}\\&=\GTenHLProbability{11/24}\approx\GTenHLProbability{0.4583}=\GTenHLProbability{45.83\%}.\end{aligned}\]
The payout for one selected three-item subset after the bonus and cap is below 12 USD with probability $\GTenHLProbability{11/24}$, approximately 45.83\%.

\textbf{(ii)} Include only $\GTenHLValue{24},\GTenHLValue{27}$ USD; include the boundary mass at 24 USD.
\[\begin{aligned}P(\GTenHLValue{X}\geq\GTenHLValue{24})&=p_X(\GTenHLValue{24})+p_X(\GTenHLValue{27})\\&=\GTenHLProbability{7/24}+\GTenHLProbability{3/40}\\&=\GTenHLProbability{11/30}\approx\GTenHLProbability{0.3667}=\GTenHLProbability{36.67\%}.\end{aligned}\]
The payout for one selected three-item subset after the bonus and cap is at least 24 USD with probability $\GTenHLProbability{11/30}$, approximately 36.67\%.



\par\smallskip
\noindent\begin{minipage}{\linewidth}\centering
\textbf{Proposed partial credit.}\par\smallskip
\begin{tabular}{@{}l r@{}}\toprule
\CreditDescription{\textbf{Assessed step or criterion}} & \textbf{Points}\\\midrule
\CreditDescription{(i): selecting exactly the included PMF masses with the correct boundary} & 0.375\\
\CreditDescription{(i): summing them correctly} & 0.1875\\
\CreditDescription{(i): the contextual interpretation} & 0.1875\\
\CreditDescription{(ii): selecting exactly the included PMF masses with the correct boundary} & 0.375\\
\CreditDescription{(ii): summing them correctly} & 0.1875\\
\CreditDescription{(ii): the contextual interpretation} & 0.1875\\
\midrule
\CreditDescription{\textbf{Total}} & \textbf{1.5}\\
\bottomrule\end{tabular}
\end{minipage}\par\smallskip

Each part totals 0.75 points.

\par\vspace{3mm}\secInicio[][sec:q6]

\textbf{Scoring reconciliation.} Question 6 gives $1.5+3+1.5=6$ points for C2. C2 gives $2+2+4+6=14$ points; its independent passing threshold is 4 points. C1 gives $3+3=6$ points; its independent passing threshold is also 4 points. The overall available total is $6+14=20$ points, with no transfer between criteria.

\end{document}
